% Options for packages loaded elsewhere
\PassOptionsToPackage{unicode}{hyperref}
\PassOptionsToPackage{hyphens}{url}
\PassOptionsToPackage{dvipsnames,svgnames,x11names}{xcolor}
\documentclass[
  11pt,
]{article}
\usepackage{xcolor}
\usepackage[margin=2.2cm]{geometry}
\usepackage{amsmath,amssymb}
\setcounter{secnumdepth}{-\maxdimen} % remove section numbering
\usepackage{iftex}
\ifPDFTeX
  \usepackage[T1]{fontenc}
  \usepackage[utf8]{inputenc}
  \usepackage{textcomp} % provide euro and other symbols
\else % if luatex or xetex
  \usepackage{unicode-math} % this also loads fontspec
  \defaultfontfeatures{Scale=MatchLowercase}
  \defaultfontfeatures[\rmfamily]{Ligatures=TeX,Scale=1}
\fi
\usepackage{lmodern}
\ifPDFTeX\else
  % xetex/luatex font selection
\fi
% Use upquote if available, for straight quotes in verbatim environments
\IfFileExists{upquote.sty}{\usepackage{upquote}}{}
\IfFileExists{microtype.sty}{% use microtype if available
  \usepackage[]{microtype}
  \UseMicrotypeSet[protrusion]{basicmath} % disable protrusion for tt fonts
}{}
\makeatletter
\@ifundefined{KOMAClassName}{% if non-KOMA class
  \IfFileExists{parskip.sty}{%
    \usepackage{parskip}
  }{% else
    \setlength{\parindent}{0pt}
    \setlength{\parskip}{6pt plus 2pt minus 1pt}}
}{% if KOMA class
  \KOMAoptions{parskip=half}}
\makeatother
\usepackage{color}
\usepackage{fancyvrb}
\newcommand{\VerbBar}{|}
\newcommand{\VERB}{\Verb[commandchars=\\\{\}]}
\DefineVerbatimEnvironment{Highlighting}{Verbatim}{commandchars=\\\{\}}
% Add ',fontsize=\small' for more characters per line
\usepackage{framed}
\definecolor{shadecolor}{RGB}{248,248,248}
\newenvironment{Shaded}{\begin{snugshade}}{\end{snugshade}}
\newcommand{\AlertTok}[1]{\textcolor[rgb]{0.94,0.16,0.16}{#1}}
\newcommand{\AnnotationTok}[1]{\textcolor[rgb]{0.56,0.35,0.01}{\textbf{\textit{#1}}}}
\newcommand{\AttributeTok}[1]{\textcolor[rgb]{0.13,0.29,0.53}{#1}}
\newcommand{\BaseNTok}[1]{\textcolor[rgb]{0.00,0.00,0.81}{#1}}
\newcommand{\BuiltInTok}[1]{#1}
\newcommand{\CharTok}[1]{\textcolor[rgb]{0.31,0.60,0.02}{#1}}
\newcommand{\CommentTok}[1]{\textcolor[rgb]{0.56,0.35,0.01}{\textit{#1}}}
\newcommand{\CommentVarTok}[1]{\textcolor[rgb]{0.56,0.35,0.01}{\textbf{\textit{#1}}}}
\newcommand{\ConstantTok}[1]{\textcolor[rgb]{0.56,0.35,0.01}{#1}}
\newcommand{\ControlFlowTok}[1]{\textcolor[rgb]{0.13,0.29,0.53}{\textbf{#1}}}
\newcommand{\DataTypeTok}[1]{\textcolor[rgb]{0.13,0.29,0.53}{#1}}
\newcommand{\DecValTok}[1]{\textcolor[rgb]{0.00,0.00,0.81}{#1}}
\newcommand{\DocumentationTok}[1]{\textcolor[rgb]{0.56,0.35,0.01}{\textbf{\textit{#1}}}}
\newcommand{\ErrorTok}[1]{\textcolor[rgb]{0.64,0.00,0.00}{\textbf{#1}}}
\newcommand{\ExtensionTok}[1]{#1}
\newcommand{\FloatTok}[1]{\textcolor[rgb]{0.00,0.00,0.81}{#1}}
\newcommand{\FunctionTok}[1]{\textcolor[rgb]{0.13,0.29,0.53}{\textbf{#1}}}
\newcommand{\ImportTok}[1]{#1}
\newcommand{\InformationTok}[1]{\textcolor[rgb]{0.56,0.35,0.01}{\textbf{\textit{#1}}}}
\newcommand{\KeywordTok}[1]{\textcolor[rgb]{0.13,0.29,0.53}{\textbf{#1}}}
\newcommand{\NormalTok}[1]{#1}
\newcommand{\OperatorTok}[1]{\textcolor[rgb]{0.81,0.36,0.00}{\textbf{#1}}}
\newcommand{\OtherTok}[1]{\textcolor[rgb]{0.56,0.35,0.01}{#1}}
\newcommand{\PreprocessorTok}[1]{\textcolor[rgb]{0.56,0.35,0.01}{\textit{#1}}}
\newcommand{\RegionMarkerTok}[1]{#1}
\newcommand{\SpecialCharTok}[1]{\textcolor[rgb]{0.81,0.36,0.00}{\textbf{#1}}}
\newcommand{\SpecialStringTok}[1]{\textcolor[rgb]{0.31,0.60,0.02}{#1}}
\newcommand{\StringTok}[1]{\textcolor[rgb]{0.31,0.60,0.02}{#1}}
\newcommand{\VariableTok}[1]{\textcolor[rgb]{0.00,0.00,0.00}{#1}}
\newcommand{\VerbatimStringTok}[1]{\textcolor[rgb]{0.31,0.60,0.02}{#1}}
\newcommand{\WarningTok}[1]{\textcolor[rgb]{0.56,0.35,0.01}{\textbf{\textit{#1}}}}
\usepackage{longtable,booktabs,array}
\newcounter{none} % for unnumbered tables
\usepackage{calc} % for calculating minipage widths
% Correct order of tables after \paragraph or \subparagraph
\usepackage{etoolbox}
\makeatletter
\patchcmd\longtable{\par}{\if@noskipsec\mbox{}\fi\par}{}{}
\makeatother
% Allow footnotes in longtable head/foot
\IfFileExists{footnotehyper.sty}{\usepackage{footnotehyper}}{\usepackage{footnote}}
\makesavenoteenv{longtable}
\setlength{\emergencystretch}{3em} % prevent overfull lines
\providecommand{\tightlist}{%
  \setlength{\itemsep}{0pt}\setlength{\parskip}{0pt}}

\usepackage{fontspec}
% Symbols missing from the text fonts (arrows, maths signs) come from these fallbacks.
\directlua{luaotfload.add_fallback("symbols", {"STIX Two Math:mode=harf;", "DejaVu Sans:mode=harf;"})}
\setmainfont{STIX Two Text}[RawFeature={fallback=symbols}]
\setmonofont{DejaVu Sans Mono}[Scale=0.85, RawFeature={fallback=symbols}]
\usepackage{unicode-math}
\setmathfont{STIX Two Math}
\usepackage{tikz}
\usetikzlibrary{graphs, graphdrawing, quotes, arrows.meta}
\usegdlibrary{layered}
\usepackage{adjustbox}
\usepackage{float}
\usepackage{fvextra}
\fvset{breaklines=true, breakanywhere=true}
\RecustomVerbatimEnvironment{verbatim}{Verbatim}{breaklines=true, breakanywhere=true}
\renewcommand{\arraystretch}{1.2}
% Let identifiers like learn_temperature=True break after underscores (e.g. in table cells).
\let\oldunderscore\_
\renewcommand{\_}{\oldunderscore\allowbreak}
\usepackage{bookmark}
\IfFileExists{xurl.sty}{\usepackage{xurl}}{} % add URL line breaks if available
\urlstyle{same}
\hypersetup{
  colorlinks=true,
  linkcolor={Maroon},
  filecolor={Maroon},
  citecolor={Blue},
  urlcolor={Blue},
  pdfcreator={LaTeX via pandoc}}

\author{}
\date{}

\begin{document}

{
\setcounter{tocdepth}{3}
\tableofcontents
}
\section{Architecture deep dive: the
model, the recipe, the
evidence}\label{architecture-deep-dive-the-model-the-recipe-the-evidence}

This is the reference for teaching the
project. It goes layer by layer through
the model, derives the loss, explains
the evaluation, lists every experiment
with its numbers, and maps each idea to
the paper it comes from.

\begin{itemize}
\tightlist
\item
  Start with
  \href{how_it_works.pdf}{how\_it\_works.md}
  for the overview (task, data, big
  picture). This document assumes it and
  goes deeper.
\item
  Code references are to commit
  \texttt{0e9ff1f} (2026-10-20) on
  branch \texttt{ws-m1-infra} (§9:
  \texttt{d3ed4ea}). Function and class
  names are stable; line numbers drift,
  so search for the name.
\item
  Every number below comes from a file
  in the repo (\texttt{RESULTS.md},
  \texttt{logs/},
  \texttt{submissions/*/run\_config*.json})
  or from a script run on the model.
  Anything that couldn\textquotesingle t
  be checked is marked
  \textbf{(unverified)}.
\end{itemize}

Contents

\begin{enumerate}
\def\labelenumi{\arabic{enumi}.}
\tightlist
\item
  \hyperref[1-the-encoder-layer-by-layer]{The
  encoder, layer by layer}
\item
  \hyperref[2-where-the-encoder-comes-from]{Where
  the encoder comes from (lineage and
  history)}
\item
  \hyperref[3-the-loss-and-the-training-recipe]{The
  loss and the training recipe}
\item
  \hyperref[4-evaluation-in-detail]{Evaluation
  in detail}
\item
  \hyperref[5-every-experiment-and-its-result]{Every
  experiment and its result}
\item
  \hyperref[6-submission-and-inference-path]{Submission
  and inference path}
\item
  \hyperref[7-paper-map-and-bibliography]{Paper
  map and bibliography}
\item
  \hyperref[8-teaching-kit]{Teaching
  kit}
\item
  \hyperref[9-next-new-people-and-new-headsets-in-progress]{Next:
  new people and new headsets (in
  progress)}
\item
  \hyperref[10-unseen-participants-and-alljoined-results-and-new-tools-oct-21]{Unseen
  participants and Alljoined: results
  and new tools (Oct 21)}
\end{enumerate}

\begin{center}\rule{0.5\linewidth}{0.5pt}\end{center}

\subsection{1. The encoder, layer by
layer}\label{the-encoder-layer-by-layer}

\texttt{EEGToImageEncoder} in
\texttt{submission/submission.py}. One
input window is a matrix of \textbf{63
electrodes × 120 time samples} (1 s at
120 Hz, from −0.2 s to +0.8 s around
image onset). The output is one
\textbf{unit-length 1536-d vector} in
DINOv2-giant\textquotesingle s embedding
space.

\subsubsection{1.1 The whole network at
a
glance}\label{the-whole-network-at-a-glance}

\begin{figure}[H]\centering
\begin{adjustbox}{max width=\linewidth, max totalheight=0.8\textheight}
\begin{tikzpicture}[>={Stealth[length=2mm]}, box/.style={draw, rounded corners=2pt, align=center, font=\footnotesize, inner sep=4pt, fill=black!4}, every edge quotes/.style={font=\scriptsize, fill=white, inner sep=1pt, align=center, auto}]
\graph [layered layout, grow=down, level sep=16mm, sibling sep=32mm, nodes={box}, edges={semithick}] {
  IN [as={EEG window}];
  U [as={unsqueeze: treat it as a 1-channel image}];
  TC [as={Temporal conv 1→32, kernel (1, 31), pad 15 \\ + BatchNorm}];
  SC [as={Spatial conv 32→64, kernel (63, 1) \\ + BatchNorm + ELU}];
  P1 [as={AvgPool (1, 4) + Dropout 0.25}];
  DW [as={Depthwise conv 64→64, kernel (1, 15), groups 64}];
  PW [as={Pointwise conv 64→128, kernel (1, 1) \\ + BatchNorm + ELU}];
  P2 [as={AvgPool (1, 4) + Dropout 0.25}];
  FL [as={Flatten}];
  H1 [as={Linear 896→512 + BatchNorm + GELU + Dropout}];
  H2 [as={Linear 512→1536}];
  N [as={L2 normalise}];
  OUT [as={Predicted image embedding}];
  IN ->["{(B, 63, 120)}"] U;
  U ->["{(B, 1, 63, 120)}"] TC;
  TC ->["{(B, 32, 63, 120)}"] SC;
  SC ->["{(B, 64, 1, 120)}"] P1;
  P1 ->["{(B, 64, 1, 30)}"] DW;
  DW ->["{(B, 64, 1, 30)}"] PW;
  PW ->["{(B, 128, 1, 30)}"] P2;
  P2 ->["{(B, 128, 1, 7)}"] FL;
  FL ->["{(B, 896)}"] H1;
  H1 ->["{(B, 512)}"] H2;
  H2 ->["{(B, 1536)}"] N;
  N ->["{(B, 1536), length 1}"] OUT;
};
\end{tikzpicture}
\end{adjustbox}
\end{figure}

\subsubsection{1.2 Layer table
(computed, not
estimated)}\label{layer-table-computed-not-estimated}

Shapes and parameter counts were
produced by running the class on a
\texttt{(2,\ 63,\ 120)} input on CPU
(script in exercise 1, §8.3).

{\def\LTcaptype{none} % do not increment counter
\begin{longtable}[]{@{}
  >{\raggedright\arraybackslash}p{(\linewidth - 6\tabcolsep) * \real{0.0741}}
  >{\raggedright\arraybackslash}p{(\linewidth - 6\tabcolsep) * \real{0.5926}}
  >{\raggedright\arraybackslash}p{(\linewidth - 6\tabcolsep) * \real{0.2346}}
  >{\raggedright\arraybackslash}p{(\linewidth - 6\tabcolsep) * \real{0.0988}}@{}}
\toprule\noalign{}
\begin{minipage}[b]{\linewidth}\raggedright
\#
\end{minipage} &
\begin{minipage}[b]{\linewidth}\raggedright
Module (exact PyTorch)
\end{minipage} &
\begin{minipage}[b]{\linewidth}\raggedright
Shape in → out
\end{minipage} &
\begin{minipage}[b]{\linewidth}\raggedright
Parameters
\end{minipage} \\
\midrule\noalign{}
\endhead
\bottomrule\noalign{}
\endlastfoot
1 &
\texttt{Conv2d(1,\ 32,\ kernel\_size=(1,\ 31),\ padding=(0,\ 15),\ bias=False)}
& (B,1,63,120) → (B,32,63,120) & 992 \\
2 & \texttt{BatchNorm2d(32)} & same &
64 \\
3 &
\texttt{Conv2d(32,\ 64,\ kernel\_size=(63,\ 1),\ groups=1,\ bias=False)}
& (B,32,63,120) → (B,64,1,120) &
\textbf{129,024} \\
4 & \texttt{BatchNorm2d(64)} ·
\texttt{ELU()} & same & 128 \\
5 &
\texttt{AvgPool2d((1,\ 4),\ stride=(1,\ 4))}
· \texttt{Dropout(0.25)} & (B,64,1,120)
→ (B,64,1,30) & 0 \\
6 &
\texttt{Conv2d(64,\ 64,\ kernel\_size=(1,\ 15),\ padding=(0,\ 7),\ groups=64,\ bias=False)}
& (B,64,1,30) → same & 960 \\
7 &
\texttt{Conv2d(64,\ 128,\ kernel\_size=(1,\ 1),\ bias=False)}
& (B,64,1,30) → (B,128,1,30) & 8,192 \\
8 & \texttt{BatchNorm2d(128)} ·
\texttt{ELU()} & same & 256 \\
9 &
\texttt{AvgPool2d((1,\ 4),\ stride=(1,\ 4))}
· \texttt{Dropout(0.25)} & (B,128,1,30)
→ (B,128,1,7) & 0 \\
10 & \texttt{flatten} & (B,128,1,7) →
(B,896) & 0 \\
11 & \texttt{Linear(896,\ 512)} &
(B,896) → (B,512) & \textbf{459,264} \\
12 & \texttt{BatchNorm1d(512)} ·
\texttt{GELU()} · \texttt{Dropout(0.25)}
& same & 1,024 \\
13 & \texttt{Linear(512,\ 1536)} &
(B,512) → (B,1536) & \textbf{787,968} \\
14 & \texttt{F.normalize(p=2,\ dim=-1)}
& same & 0 \\
& \textbf{Total} & &
\textbf{1,387,872} \\
\end{longtable}
}

\textbf{Where the parameters are:} the
convolutional "feature extractor"
(layers 1--9) has 139,616 parameters
(10\%); the projection head (11--13) has
1,248,256 (90\%). The network is mostly
a mapping from 896 EEG features to
DINOv2 space. The single biggest conv
layer is the spatial one (129k), because
every output filter sees all 32 temporal
maps on all 63 electrodes.

Size variants from \#27
(\texttt{temporal\_filters},
\texttt{spatial\_filters},
\texttt{sep\_filters},
\texttt{hidden\_dim},
\texttt{head\_blocks}), also computed:
wider convs 64/128/256 → 2.26M; head
1024 + 1 residual block → 5.00M; both →
6.33M.

\subsubsection{1.3 What each layer does,
and
why}\label{what-each-layer-does-and-why}

\textbf{Input reshaping.} EEG is not an
image: one axis is space (electrodes),
the other is time. Treating it as a
1-channel image \texttt{(1,\ 63,\ 120)}
lets us use 2-D convolutions whose
kernels are 1-D in one axis:
\texttt{(1,\ k)} kernels act
\textbf{along time only},
\texttt{(63,\ 1)} kernels \textbf{across
electrodes only}. This "factorised"
spatio-temporal design is the core idea
of ShallowConvNet / DeepConvNet
(Schirrmeister et al., 2017) and EEGNet
(Lawhern et al., 2018).

\textbf{1--2. Temporal convolution (32
filters × 31 taps) + BatchNorm.}

\begin{itemize}
\tightlist
\item
  Each filter is a \textbf{learned FIR
  filter} applied identically to every
  electrode (the kernel is
  \texttt{(1,\ 31)}, so it slides along
  time and is shared across rows).
\item
  31 taps at 120 Hz span \textbf{258
  ms}. A filter of that length can
  isolate frequency bands about 120/31 ≈
  4 Hz wide anywhere between 0 and the
  60 Hz Nyquist limit: theta (4--8 Hz),
  alpha (8--13 Hz), beta, or act as a
  matched filter for the shape of an
  evoked response.
\item
  \texttt{padding=15} keeps the 120
  samples ("same" convolution). The
  first and last 15 samples are computed
  against zero padding, so the window
  edges are slightly less reliable.
\item
  \texttt{bias=False} because the
  BatchNorm right after has its own
  shift.
\item
  \textbf{BatchNorm2d} (Ioffe \&
  Szegedy, 2015) standardises each
  filter\textquotesingle s output over
  the batch, electrodes and time, which
  keeps training stable at a high
  learning rate.
\item
  Lineage: EEGNet uses 8 filters of half
  the sampling rate (64 taps at 128 Hz ≈
  500 ms); ShallowConvNet uses 40
  filters of 25 taps at 250 Hz (100 ms).
  Ours sits in between, with more
  filters.
\end{itemize}

\textbf{3--4. Spatial convolution (64
filters spanning all 63 electrodes) +
BatchNorm + ELU.}

\begin{itemize}
\tightlist
\item
  Kernel \texttt{(63,\ 1)}: each output
  is a \textbf{weighted sum over all
  electrodes} at one time point. This is
  a learned \textbf{spatial filter}, the
  deep-learning version of classical
  techniques like Common Spatial
  Patterns (CSP), beamformers, or
  picking an ICA component: it can, for
  example, weight the occipital
  electrodes positively and subtract a
  frontal reference.
\item
  \texttt{groups=1} means each of the 64
  spatial filters mixes \textbf{all 32
  temporal maps}: it combines frequency
  bands and electrodes jointly (32 × 63
  = 2,016 weights per filter). EEGNet
  instead uses a \textbf{depthwise}
  spatial conv (each spatial filter sees
  only one temporal map,
  \texttt{D\ =\ 2} spatial filters per
  temporal filter, 63 weights each, with
  a max-norm constraint). Ours has many
  more parameters here (129k vs
  EEGNet\textquotesingle s 1,008 at the
  same input size). \emph{Note:} the
  original scaffold\textquotesingle s
  comment calls this layer "depthwise"
  but the code has always used
  \texttt{groups=1} (commit
  \texttt{db3fb74}); the code is what
  ran in every experiment.
\item
  The kernel height is the channel
  count, so \textbf{the layer is tied to
  63 electrodes in a fixed order}.
  That\textquotesingle s why a
  THINGS-EEG2 model
  can\textquotesingle t run on the final
  phase\textquotesingle s 32 Emotiv
  channels without retraining (see issue
  \#21 and the planned position-based
  spatial attention).
\item
  \textbf{ELU} (exponential linear unit,
  Clevert et al., 2016) is smooth for
  negative inputs; EEGNet and
  ShallowConvNet-family models use it.
\end{itemize}

\textbf{5. Average pooling (×4) +
Dropout.}

\begin{itemize}
\tightlist
\item
  120 → 30 time steps: an average over 4
  samples is a crude \textbf{low-pass
  filter followed by downsampling} to 30
  Hz (33 ms per step). Evoked visual
  responses are slow (most power below
  \textasciitilde20 Hz), so little is
  lost and later layers become cheaper
  and see longer stretches of time.
\item
  \textbf{Dropout 0.25} randomly zeroes
  25\% of activations during training
  (Srivastava et al., 2014).
  EEGNet\textquotesingle s paper
  recommends 0.25 for cross-subject
  training (0.5 within-subject)
  \textbf{(from memory of the EEGNet
  paper; braindecode\textquotesingle s
  default is 0.25)}.
\end{itemize}

\textbf{6--8. Depthwise-separable
convolution + BatchNorm + ELU.}

\begin{itemize}
\tightlist
\item
  \textbf{Depthwise} (layer 6): each of
  the 64 maps gets its \textbf{own}
  15-tap temporal filter
  (\texttt{groups=64}, 960 weights). At
  30 Hz, 15 taps span \textbf{500 ms}:
  it summarises how each spatial
  component evolves over half a second.
\item
  \textbf{Pointwise} (layer 7): a 1×1
  conv that mixes the 64 maps into 128
  (8,192 weights).
\item
  Splitting "filter in time" from "mix
  the maps" is a
  \textbf{depthwise-separable
  convolution} (Chollet, Xception, CVPR
  2017), adopted by EEGNet to use far
  fewer parameters than a full conv (a
  full 64→128 conv with 15 taps would
  need 122,880 weights instead of
  9,152).
\end{itemize}

\textbf{9. Average pooling (×4) +
Dropout.} 30 → 7 time steps (30/4 = 7.5,
floored: the last two 30 Hz steps are
not pooled, though their information
still reaches step 6 through the
depthwise kernel). Each final step
represents 16 input samples =
\textbf{133 ms}.

\textbf{10. Flatten → 896 features.} 128
maps × 7 time steps. Flattening (rather
than averaging over time) keeps
\textbf{when} each feature occurred: the
next layer has separate weights for
"component 12 at 100--233 ms" and
"component 12 at 500--633 ms". That
matters because the timing of visual
responses carries information. It is
also our best explanation for why random
time-shift augmentation hurt (R4, §5).

\textbf{Receptive field of the final
features.} Following the layers back to
the input: temporal conv (31) → pool
(stride 4) → depthwise conv (15 at
stride 4) → pool (stride 16). Each of
the 7 final time steps sees \textbf{102
input samples = 850 ms}, and consecutive
steps are \textbf{133 ms apart}. Step 0
covers roughly −200 → +283 ms around
image onset; step 6 roughly +242 → +792
ms. So every final feature integrates
most of the window, with a different
centre of gravity: early steps emphasise
the first visual responses
(\textasciitilde100--200 ms), late steps
the later, more semantic activity.

\textbf{11--13. Projection head.}
\texttt{Linear(896→512)\ →\ BatchNorm1d\ →\ GELU\ →\ Dropout\ →\ Linear(512→1536)}.

\begin{itemize}
\tightlist
\item
  A two-layer MLP that maps EEG features
  into DINOv2\textquotesingle s 1536-d
  space. Using an MLP "projector"
  between the encoder and the
  contrastive loss is the standard
  recipe in contrastive learning (CLIP
  uses linear projections; SimCLR showed
  a non-linear head helps).
\item
  \textbf{GELU} (Hendrycks \& Gimpel,
  2016) is a smooth ReLU used in
  transformers.
\item
  It\textquotesingle s 90\% of the
  model. The NeuralBench EEGNet baseline
  instead reads 48 features (16 maps × 3
  time steps) straight into 1536 outputs
  with a single linear layer; that
  "bigger head" is one of the two
  changes \texttt{RESULTS.md} credits
  for our jump from 0.30 to 0.54 (§2.2).
\item
  Optional (\#27): \texttt{head\_blocks}
  adds residual blocks
  \texttt{x\ +\ Dropout(Linear(GELU(x)))}
  on the 1536-d output followed by
  \texttt{LayerNorm}, as in
  ATM\textquotesingle s projector (Li et
  al., 2024). Off by default; R13/R14 in
  the queue test it.
\end{itemize}

\textbf{14. L2 normalisation.} Divides
the output by its length so it lies on
the unit sphere. The metric ranks by
\textbf{cosine similarity}, which only
depends on direction, and the loss
normalises the targets too, so
predictions and targets are compared
purely by angle.

\begin{center}\rule{0.5\linewidth}{0.5pt}\end{center}

\subsection{2. Where the encoder comes
from}\label{where-the-encoder-comes-from}

\subsubsection{2.1 Family
tree}\label{family-tree}

\begin{figure}[H]\centering
\begin{adjustbox}{max width=\linewidth, max totalheight=0.8\textheight}
\begin{tikzpicture}[>={Stealth[length=2mm]}, box/.style={draw, rounded corners=2pt, align=center, font=\footnotesize, inner sep=4pt, fill=black!4}, every edge quotes/.style={font=\scriptsize, fill=white, inner sep=1pt, align=center, auto}]
\graph [layered layout, grow=down, level sep=16mm, sibling sep=32mm, nodes={box}, edges={semithick}] {
  FBCSP [as={Filter-bank CSP \\ (classical BCI: \\ band-pass, then spatial filters)}];
  SH [as={ShallowConvNet / DeepConvNet \\ Schirrmeister et al. 2017 \\ temporal conv → spatial conv}];
  EEGNET [as={EEGNet \\ Lawhern et al. 2018 \\ + depthwise spatial conv \\ + separable conv, ELU, dropout}];
  XC [as={Xception \\ Chollet 2017 \\ depthwise-separable conv}];
  NB [as={NeuralBench EEGNet baseline \\ 77k params, linear readout \\ our 0.29}];
  OURS [as={Our EEGToImageEncoder \\ full spatial conv, wider, \\ 512-d MLP head, 1.39M params \\ our 0.54 → 0.74}];
  CLIP [as={CLIP \\ Radford et al. 2021 \\ contrastive loss, projector, \\ learnable temperature}];
  NICE [as={NICE \\ Song et al. 2024 \\ ShallowNet-style patch encoder \\ + residual projector}];
  ATM [as={ATM \\ Li et al. 2024 \\ channel attention + conv \\ + residual projector}];
  FBCSP ->[] SH;
  SH ->[] EEGNET;
  XC ->[] EEGNET;
  EEGNET ->[] NB;
  EEGNET ->[] OURS;
  CLIP ->[] OURS;
  SH ->[] NICE;
  NICE ->[] ATM;
  ATM ->[dashed, "{head\_blocks option, issue 27}"] OURS;
};
\end{tikzpicture}
\end{adjustbox}
\end{figure}

\begin{itemize}
\tightlist
\item
  \textbf{Filter-bank CSP} is the
  classical pipeline for EEG decoding:
  band-pass filter the signal into
  several frequency bands, learn spatial
  filters per band, then classify the
  log-power.
\item
  \textbf{ShallowConvNet} (Schirrmeister
  et al., 2017) turned that into a
  network: temporal conv (40 filters ×
  25 taps) → spatial conv (40) →
  squaring → mean pooling (75, stride
  15) → log → classifier, i.e. a
  learnable filter-bank CSP.
  \textbf{DeepConvNet} stacks four
  conv-pool blocks. (Defaults verified
  in braindecode\textquotesingle s
  \texttt{ShallowFBCSPNet}: 40 / 25 / 40
  / 75 / 15.)
\item
  \textbf{EEGNet} (Lawhern et al., 2018)
  made it compact: temporal conv (F1 = 8
  filters, 64 taps) → depthwise spatial
  conv (D = 2 per temporal filter,
  max-norm constrained) → BatchNorm,
  ELU, average pooling, dropout →
  separable conv (F2 = 16) → pooling →
  classifier. (Defaults verified in
  braindecode\textquotesingle s
  \texttt{EEGNet}: F1 = 8, D = 2, F2 =
  F1·D, kernel 64, depthwise kernel 16,
  dropout 0.25.)
\item
  \textbf{Our encoder} keeps
  EEGNet\textquotesingle s block order
  and its separable conv, but uses 4×
  more temporal filters (32), a
  \textbf{full} (non-depthwise,
  unconstrained) spatial conv with 64
  filters, 128 separable maps, a shorter
  temporal kernel (31 taps at 120 Hz),
  and replaces the classifier with a
  2-layer \textbf{projection head} into
  DINOv2 space trained with a CLIP loss.
\item
  \textbf{NICE} (Song et al., ICLR 2024)
  and \textbf{ATM} (Li et al., NeurIPS
  2024) are the THINGS-EEG2 state of the
  art that issue \#12 targets. From the
  papers and their public code, roughly
  \textbf{(unverified in detail; check
  the papers\textquotesingle{}
  architecture figures)}:
  NICE\textquotesingle s encoder is a
  ShallowNet-style temporal conv →
  average pooling → spatial conv
  producing a short sequence of patch
  embeddings, then a projector
  \texttt{Linear\ →\ residual(GELU,\ Linear,\ Dropout)\ →\ LayerNorm};
  NICE-GA adds graph attention over
  electrodes. ATM adds a channel-wise
  transformer embedding (with a subject
  token) before a similar conv block and
  residual projector. Our
  \texttt{head\_blocks} option is that
  residual projector.
\end{itemize}

\subsubsection{2.2 The true history
(from
git)}\label{the-true-history-from-git}

{\def\LTcaptype{none} % do not increment counter
\begin{longtable}[]{@{}
  >{\raggedright\arraybackslash}p{(\linewidth - 4\tabcolsep) * \real{0.1341}}
  >{\raggedright\arraybackslash}p{(\linewidth - 4\tabcolsep) * \real{0.2561}}
  >{\raggedright\arraybackslash}p{(\linewidth - 4\tabcolsep) * \real{0.6098}}@{}}
\toprule\noalign{}
\begin{minipage}[b]{\linewidth}\raggedright
Date
\end{minipage} &
\begin{minipage}[b]{\linewidth}\raggedright
Commit
\end{minipage} &
\begin{minipage}[b]{\linewidth}\raggedright
What happened
\end{minipage} \\
\midrule\noalign{}
\endhead
\bottomrule\noalign{}
\endlastfoot
2026-09-11 & \texttt{db3fb74} & Initial
scaffold.
\texttt{src/models/eeg\_encoder.py}
defines the architecture as
\textbf{learning TODOs} "based on EEGNet
principles": temporal conv (1, 31) × 32
→ spatial conv (n\_chans, 1) × 64 →
separable conv → 512-d projection head →
1536, L2-normalised. The layer sizes we
still use were fixed here. \\
2026-09-13 & \texttt{7d7dc1e} & TODOs
completed ("complete TODOs for
setup"). \\
2026-10-17/18 & \texttt{6896932},
\texttt{bf5bf17} &
\texttt{submission/submission.py} made
self-contained (torch only) and fixed
for the Codabench contract; state\_dict
keys match
\texttt{src/models/eeg\_encoder.py}. \\
2026-10-18 & \texttt{7b439e0} &
\textbf{First real upload:
NeuralBench\textquotesingle s own
EEGNet} (official template), 0.29. \\
2026-10-19 & \texttt{da20e91},
\texttt{c652e3a} & Our encoder trained
on THINGS-EEG2 with InfoNCE:
\textbf{0.54} (held-out 42.15\%). Seeded
training and \texttt{weights\_*.pt}
ensembling added. \\
2026-10-19 & \texttt{5facb9b} & Training
options for the round-1 experiments
(\texttt{learn\_temperature},
\texttt{dedup}, \texttt{symmetric},
augmentation) and \texttt{-\/-set}
overrides. \\
2026-10-19 & \texttt{1a8fc30},
\texttt{41b152e}, \texttt{2fb5a51} &
Learnable temperature (+10 held-out),
\texttt{-\/-use-val} (+3.5). \\
2026-10-20 & \texttt{e34d3f1} & Final
3-seed ensemble on all concepts: 73.80\%
test, \textbf{0.74} on Codabench. \\
2026-10-20 & \texttt{18992e0},
\texttt{424c9b5} & Per-epoch held-out
scoring + early stopping (\#15);
repetition averaging (\#17). \\
2026-10-20 & \texttt{e893e30},
\texttt{c5881bd}, \texttt{48925dd} &
Schedule / weight decay / dropout
settings (\#18); encoder size settings,
sizes read back from weights (\#27);
\texttt{min\_temperature}. \\
2026-10-20 & \texttt{68d2b65},
\texttt{0e9ff1f} & Window cache (\#14);
per-epoch checkpoints and
\texttt{-\/-resume}; pause/resume. \\
\end{longtable}
}

\textbf{Why our encoder beat the
NeuralBench baseline by 24 points}
(\texttt{RESULTS.md}: "bigger head, 10×
higher lr"). The
baseline\textquotesingle s configuration
(\texttt{submissions/neuralbench\_eegnet\_things\_eeg2\_seed34/run\_config.yaml})
versus ours:

{\def\LTcaptype{none} % do not increment counter
\begin{longtable}[]{@{}
  >{\raggedright\arraybackslash}p{(\linewidth - 4\tabcolsep) * \real{0.0761}}
  >{\raggedright\arraybackslash}p{(\linewidth - 4\tabcolsep) * \real{0.5435}}
  >{\raggedright\arraybackslash}p{(\linewidth - 4\tabcolsep) * \real{0.3804}}@{}}
\toprule\noalign{}
\begin{minipage}[b]{\linewidth}\raggedright
\end{minipage} &
\begin{minipage}[b]{\linewidth}\raggedright
NeuralBench EEGNet (0.29--0.30)
\end{minipage} &
\begin{minipage}[b]{\linewidth}\raggedright
Our encoder (0.54)
\end{minipage} \\
\midrule\noalign{}
\endhead
\bottomrule\noalign{}
\endlastfoot
Network & braindecode EEGNet defaults
(F1 8, D 2, F2 16, kernel 64),
\textbf{77k parameters} (counted from
the shipped weights) & 1.39M
parameters \\
Readout & one linear layer: 48 features
(16 maps × 3 steps) → 1536 & 896
features → 512 → 1536 MLP \\
Optimiser & AdamW, \textbf{lr 1e-4},
weight decay 0.05, OneCycle schedule &
AdamW, \textbf{lr 1e-3}, weight decay
1e-4, constant \\
Batch & 64 & 256 (4× more negatives per
step) \\
Loss & NeuralBench \texttt{ClipLoss},
fixed temperature & our
\texttt{clip\_loss}, τ = 0.07 fixed
(learnable later) \\
Stopping & 40 epochs, patience 5 on a
validation rank metric & 20 epochs \\
\end{longtable}
}

We never ran the ablation that separates
these factors, so the split between
"bigger head", "higher lr" and "bigger
batch" is \textbf{not measured}.

\begin{center}\rule{0.5\linewidth}{0.5pt}\end{center}

\subsection{3. The loss and the training
recipe}\label{the-loss-and-the-training-recipe}

All in \texttt{Solver.fit} and
\texttt{clip\_loss} in
\texttt{submission/submission.py},
driven by
\texttt{scripts/train\_solver.py}.

\subsubsection{3.1 The contrastive
(InfoNCE / CLIP)
loss}\label{the-contrastive-infonce-clip-loss}

For a batch of \(B\) windows, let
\(p_i\) be the model\textquotesingle s
unit-length prediction for window \(i\)
and \(t_i\) the DINOv2 embedding of the
image that window shows (normalised
inside the loss). The similarity matrix
is \(s_{ij} = p_i \cdot t_j\) (cosine).
With temperature \(\tau\):

\[
\mathcal{L} = -\frac{1}{B}\sum_{i=1}^{B} \log \frac{\exp(s_{ii}/\tau)}{\sum_{j=1}^{B}\exp(s_{ij}/\tau)}
\]

Read it as a \textbf{B-way
classification per row}: "which of the B
images in this batch did EEG window
\(i\) see?" The correct answer is the
diagonal; the other \(B-1\) images in
the batch are the \textbf{negatives}
("in-batch negatives"), so no extra
sampling is needed. Minimising it pulls
\(p_i\) towards \(t_i\) and pushes it
away from the other targets, which is
exactly what top-k retrieval rewards.
This is \textbf{InfoNCE} (van den Oord
et al., 2018), popularised by
\textbf{CLIP} (Radford et al., 2021),
and used for brain decoding by Défossez
et al. (2023), NICE, ATM.

Why not regression (minimise
\(\|p_i - t_i\|^2\))? A regression model
facing noisy input predicts the
conditional mean, a blurry "average
image" that sits close to everything and
ranks poorly. \#23\textquotesingle s
post-hoc ridge maps confirmed it: fitted
by squared error, they scored
0.046--0.052 held-out (§5).

\textbf{Options tested in round 1:}

\begin{itemize}
\tightlist
\item
  \texttt{symmetric=True} adds the
  column-wise term (for each image,
  which EEG window matches?), as in
  CLIP. No gain here (R3).
\item
  \texttt{dedup=True} masks out "false
  negatives": other windows in the batch
  that show the \textbf{same} image
  (each training image has
  \textasciitilde40 windows: 10 subjects
  × 4 repetitions). With
  \textasciitilde14.5k training images
  and batch 256, a batch contains only
  about 2 such pairs on average
  \textbf{(our estimate)}, which
  explains why it made no difference
  (R1: 41.20\% vs 41.35\%).
\end{itemize}

\subsubsection{3.2 The
temperature}\label{the-temperature}

\(\tau\) divides the similarities before
the softmax. Small \(\tau\) makes the
softmax sharp, so the loss concentrates
on the hardest negatives (images that
look most like the right one).

\begin{itemize}
\tightlist
\item
  Fixed at \textbf{0.07} at first (a
  common default).
\item
  \textbf{Learnable}
  (\texttt{learn\_temperature=True}),
  parameterised as CLIP does: the
  trainable parameter is
  \(\log(1/\tau)\) (initialised to
  \(\log(1/0.07) \approx 2.66\)) and it
  is clamped at \(\log 100\), i.e.
  \textbf{\(\tau \ge 0.01\)} ("logit
  scale at most 100").
  It\textquotesingle s trained with the
  model by the same optimiser.
\item
  Effect: \textbf{+10 points held-out}
  (R2: 41.35\% → 51.75\%), the biggest
  single gain.
\item
  In R6 the learned \(\tau\) fell from
  0.027 (epoch 1) to the \textbf{0.01
  floor by epoch \textasciitilde15} and
  stayed there. R15 lowered the floor to
  0.005 (\texttt{min\_temperature}):
  \(\tau\) went to 0.0072, but held-out
  trailed R6 at every comparable epoch
  (best 0.564 vs 0.582 by epoch 31) when
  a restart stopped it at epoch 32.
  \textbf{Conclusion: no gain}, floor
  stays at 0.01.
\end{itemize}

\subsubsection{3.3
Optimisation}\label{optimisation}

{\def\LTcaptype{none} % do not increment counter
\begin{longtable}[]{@{}
  >{\raggedright\arraybackslash}p{(\linewidth - 4\tabcolsep) * \real{0.1304}}
  >{\raggedright\arraybackslash}p{(\linewidth - 4\tabcolsep) * \real{0.3261}}
  >{\raggedright\arraybackslash}p{(\linewidth - 4\tabcolsep) * \real{0.5435}}@{}}
\toprule\noalign{}
\begin{minipage}[b]{\linewidth}\raggedright
Setting
\end{minipage} &
\begin{minipage}[b]{\linewidth}\raggedright
Value
\end{minipage} &
\begin{minipage}[b]{\linewidth}\raggedright
Notes / source
\end{minipage} \\
\midrule\noalign{}
\endhead
\bottomrule\noalign{}
\endlastfoot
Optimiser & \textbf{AdamW} & Adam with
decoupled weight decay (Loshchilov \&
Hutter, ICLR 2019) \\
Learning rate & \textbf{1e-3} & 10× the
NeuralBench baseline \\
Weight decay & \textbf{1e-4} &
\texttt{weight\_decay} setting (\#18);
applied to the temperature parameter too
(negligible) \\
Batch size & \textbf{256} &
\texttt{-\/-batch-size}; more negatives
per step. R10/R11 test 512 / 1024 with
linearly scaled lr (Goyal et al.,
2017) \\
Gradient clipping & max norm
\textbf{1.0} & stabilises early
training \\
Schedule & constant (default);
\texttt{lr\_schedule=\textquotesingle{}cosine\textquotesingle{}}
+ \texttt{warmup\_epochs} & Cosine decay
from SGDR (Loshchilov \& Hutter, ICLR
2017); linear warm-up from Goyal et al.
2017. Applied per step via
\texttt{LambdaLR} \\
Epochs & 20 originally; \textbf{40} in
R6 (+3.75 held-out) & longer training
still improving at 40 \\
Dropout & 0.25 & \texttt{dropout}
setting \\
Seed & \texttt{-\/-seed} & seeds weight
init, shuffling and dropout \\
\end{longtable}
}

\subsubsection{3.4 Data-side
recipe}\label{data-side-recipe}

\begin{itemize}
\tightlist
\item
  \textbf{More data
  (\texttt{-\/-use-val}, \#16):}
  NeuralBench\textquotesingle s own
  "val" split of THINGS-EEG2 is ordinary
  training data (same concepts). Adding
  it gives +25\% windows (470,140 →
  587,574 with the held-out concepts
  removed; 668,400 for final models on
  all concepts): \textbf{+3.5 held-out}.
\item
  \textbf{Augmentation (R4):} random
  circular time shift up to ±6 samples
  (±50 ms) and 10\% channel dropout. It
  \textbf{hurt badly} (49.85\% →
  36.80\%). Our interpretation
  \textbf{(hypothesis, not tested)}: the
  flattened readout uses absolute timing
  (§1.3, layer 10), the circular roll
  wraps the end of the window onto the
  start, and 50 ms is large relative to
  visual response latencies.
\item
  \textbf{Repetition averaging
  (\texttt{-\/-avg-reps\ K}, \#17,
  Gifford et al. averaged repetitions
  too):} each training item becomes the
  mean of K repetitions of the same
  (subject, image), which cuts
  independent noise by about
  \(\sqrt{K}\). With
  \texttt{-\/-avg-reps-random} K is
  drawn from 1..K per item, mixing clean
  and noisy examples (test inputs are
  single trials). There are 145,400
  (subject, image) groups, so an "epoch"
  is \textbf{4× shorter} than with
  single windows. \textbf{Caveat:} early
  stopping counts epochs, so patience 5
  is effectively 4× less patient. R7
  (random 1--4) was still rising when it
  stopped at 0.555 (R6: 0.590), but it
  got there in 31 min, where R6 needed
  \textasciitilde105 min for 0.59 and
  was at \textasciitilde0.51 after 31
  min. A fair test needs patience scaled
  to the epoch length.
\end{itemize}

\subsubsection{3.5 Training
infrastructure (no effect on results, by
design)}\label{training-infrastructure-no-effect-on-results-by-design}

\begin{itemize}
\tightlist
\item
  \textbf{Early stopping on held-out
  (\#15):} after each epoch,
  \texttt{scripts/train\_solver.py}
  scores the held-out split with the
  official objective; \texttt{fit} keeps
  the best epoch\textquotesingle s
  weights and stops after
  \texttt{patience} epochs without
  improvement (5; 8 for cosine runs).
  Caveat for teaching: picking the best
  epoch on held-out makes the reported
  held-out score slightly optimistic (a
  mild selection bias); comparisons
  between runs remain fair because every
  run gets it.
\item
  \textbf{Window cache (\#14,
  \texttt{scripts/cache\_windows.py},
  \texttt{-\/-cache}):} windows
  extracted once to memory-mapped
  \texttt{.npy} files (24 GB). Verified
  to give bit-identical windows, scores
  and training batches, with
  \textasciitilde4× faster loading and
  much less memory.
\item
  \textbf{Checkpoints and
  \texttt{-\/-resume}:} \texttt{fit}
  saves its full state (model,
  optimiser, scheduler, temperature,
  best epoch, history, RNG states) after
  every epoch; a resumed run on CPU ends
  with bit-identical weights.
  \texttt{make\ pause-training} /
  \texttt{resume-training} freeze and
  continue a run (used when the
  laptop\textquotesingle s GPU is needed
  for the screen).
\end{itemize}

\begin{center}\rule{0.5\linewidth}{0.5pt}\end{center}

\subsection{4. Evaluation in
detail}\label{evaluation-in-detail}

\subsubsection{4.1 The official
objective}\label{the-official-objective}

\texttt{external/2026-competition/tracks/image\_decoding/objective.py},
\texttt{Objective.evaluate\_result}:

\begin{figure}[H]\centering
\begin{adjustbox}{max width=\linewidth, max totalheight=0.8\textheight}
\begin{tikzpicture}[>={Stealth[length=2mm]}, box/.style={draw, rounded corners=2pt, align=center, font=\footnotesize, inner sep=4pt, fill=black!4}, every edge quotes/.style={font=\scriptsize, fill=white, inner sep=1pt, align=center, auto}]
\graph [layered layout, grow=down, level sep=9mm, sibling sep=6mm, nodes={box}, edges={semithick}] {
  L [as={Test loader: (X, y, info) batches}];
  P [as={pred = model.predict(X) \\ one 1536-d vector per window}];
  Y [as={y = target embedding of the image shown}];
  G [as={Gallery = np.unique(y): the distinct images \\ (200 for THINGS-EEG2 test)}];
  S1 [as={Pooled: cosine(pred, gallery) per window}];
  GM [as={group\_means over (subject, image): \\ average a subject's predictions for the same image}];
  S2 [as={Subject-aggregated: cosine(mean pred, gallery) \\ 2,000 queries = 10 subjects × 200 images}];
  K1 [as={top5\_acc, top1\_acc}];
  K2 [as={top5\_acc\_subject\_agg, top1\_acc\_subject\_agg}];
  L ->[] P;
  L ->[] Y;
  Y ->[] G;
  P ->[] S1;
  P ->[] GM;
  GM ->[] S2;
  S1 ->[] K1;
  S2 ->[] K2;
  G ->[] S1;
  G ->[] S2;
};
\end{tikzpicture}
\end{adjustbox}
\end{figure}

\begin{itemize}
\tightlist
\item
  \textbf{Gallery:} the unique target
  embeddings of the split, found with
  \texttt{np.unique(y\_true,\ axis=0)}.
  Repeated presentations of an image
  share an identical embedding, so they
  collapse to one candidate.
\item
  \textbf{Scores:} predictions and
  candidates are L2-normalised; the
  score is the dot product (cosine
  similarity).
\item
  \textbf{Top-k accuracy}
  (\texttt{benchmark\_utils/metrics.py},
  \texttt{topk\_accuracy}): a query is a
  hit if its true
  image\textquotesingle s rank is below
  k. Ties share their average rank.
  \textbf{Chance level}: 5/200 =
  \textbf{2.5\%} for top-5, 0.5\% for
  top-1.
\item
  \textbf{Subject-aggregated}
  (\texttt{group\_means}): average the
  predictions of all windows with the
  same (subject, image) key, then rank
  those means. On the test set each
  subject saw each image
  \textasciitilde80 times, so this
  averages \textasciitilde80 predictions
  per query: it\textquotesingle s why
  our top-5 is 0.74 subject-aggregated
  but 0.31 pooled. The sealed phase
  ranks on top-5 subject-aggregated; the
  warm-up leaderboard\textquotesingle s
  main column has matched our
  subject-aggregated top-5 on every
  upload. (The
  objective\textquotesingle s docstring
  says warm-up uses the pooled metric;
  the leaderboard columns show both, and
  README/RESULTS track
  subject-aggregated. \textbf{Worth
  confirming with the organisers.})
\end{itemize}

\subsubsection{4.2 Our held-out-concept
split}\label{our-held-out-concept-split}

\texttt{src/data/heldout.py} +
\texttt{configs/heldout\_concepts.json}
(seed 0, committed):

\begin{itemize}
\tightlist
\item
  200 of the 1,654 training
  \textbf{concepts} are removed from
  training entirely (no window of those
  concepts is trained on), so the model
  must generalise to unseen concepts,
  like the test set.
\item
  One image per held-out concept (e.g.
  \texttt{air\_conditioner\_06s}) forms
  a \textbf{200-image gallery}, the same
  size as the test gallery; its windows
  (from NeuralBench\textquotesingle s
  train and val splits) are the queries:
  8,076 windows ≈ 200 images × 10
  subjects × \textasciitilde4
  repetitions.
\item
  Training then uses the remaining 1,454
  concepts (587,574 windows with
  \texttt{-\/-use-val}).
\end{itemize}

\textbf{Why the held-out score is lower
than the test score for the same model:}
the held-out images were shown only
\textbf{\textasciitilde4 times} per
subject (training-set protocol) while
test images were shown
\textbf{\textasciitilde80 times}.
Subject aggregation averages
\textasciitilde4 predictions instead of
\textasciitilde80, so much less noise is
removed. The ranking of models is what
matters, and it has been consistent.

\subsubsection{4.3 From held-out to the
leaderboard
(empirical)}\label{from-held-out-to-the-leaderboard-empirical}

{\def\LTcaptype{none} % do not increment counter
\begin{longtable}[]{@{}
  >{\raggedright\arraybackslash}p{(\linewidth - 6\tabcolsep) * \real{0.2826}}
  >{\raggedright\arraybackslash}p{(\linewidth - 6\tabcolsep) * \real{0.2391}}
  >{\raggedright\arraybackslash}p{(\linewidth - 6\tabcolsep) * \real{0.3478}}
  >{\raggedright\arraybackslash}p{(\linewidth - 6\tabcolsep) * \real{0.1304}}@{}}
\toprule\noalign{}
\begin{minipage}[b]{\linewidth}\raggedright
Run
\end{minipage} &
\begin{minipage}[b]{\linewidth}\raggedright
Held-out top-5 SA
\end{minipage} &
\begin{minipage}[b]{\linewidth}\raggedright
Test top-5 SA (same model)
\end{minipage} &
\begin{minipage}[b]{\linewidth}\raggedright
Gap
\end{minipage} \\
\midrule\noalign{}
\endhead
\bottomrule\noalign{}
\endlastfoot
control (seed 1) & 0.4135 & 0.5450 &
+13.2 \\
R2 learnable τ & 0.5175 & 0.6525 &
+13.5 \\
R5 + val data & 0.5525 & 0.7110 &
+15.9 \\
\end{longtable}
}

Then retraining the chosen recipe on
\textbf{all} concepts with 3 seeds and
ensembling added +2.7 on test (R5 0.7110
→ 0.7380). So a held-out gain of x
points has translated into roughly x
points on the leaderboard.
R6\textquotesingle s 0.590 suggests
about \textbf{0.75--0.78} for its final
ensemble \textbf{(a forecast, not a
measurement)}.

\textbf{Why local test = Codabench:}
same data pipeline (NeuralBench pinned
to the worker\textquotesingle s version,
\texttt{neuralbench==0.4.0} etc.), same
test split, same objective code,
deterministic inference. All six uploads
so far matched local scores to two
decimals.

\begin{center}\rule{0.5\linewidth}{0.5pt}\end{center}

\subsection{5. Every experiment and its
result}\label{every-experiment-and-its-result}

Held-out = top-5 subject-aggregated on
our held-out split (seed 1 unless
noted). Test = top-5 subject-aggregated
on THINGS-EEG2 test, computed locally.
Codabench = leaderboard score. Sources:
\texttt{RESULTS.md},
\texttt{logs/overnight\_summary.md},
\texttt{logs/round2.status},
\texttt{logs/*.log},
\texttt{submissions/*/run\_config.json}.

{\def\LTcaptype{none} % do not increment counter
\begin{longtable}[]{@{}
  >{\raggedright\arraybackslash}p{(\linewidth - 12\tabcolsep) * \real{0.1117}}
  >{\raggedright\arraybackslash}p{(\linewidth - 12\tabcolsep) * \real{0.1489}}
  >{\raggedright\arraybackslash}p{(\linewidth - 12\tabcolsep) * \real{0.2287}}
  >{\raggedright\arraybackslash}p{(\linewidth - 12\tabcolsep) * \real{0.1330}}
  >{\raggedright\arraybackslash}p{(\linewidth - 12\tabcolsep) * \real{0.0851}}
  >{\raggedright\arraybackslash}p{(\linewidth - 12\tabcolsep) * \real{0.0479}}
  >{\raggedright\arraybackslash}p{(\linewidth - 12\tabcolsep) * \real{0.2447}}@{}}
\toprule\noalign{}
\begin{minipage}[b]{\linewidth}\raggedright
Run
\end{minipage} &
\begin{minipage}[b]{\linewidth}\raggedright
Hypothesis
\end{minipage} &
\begin{minipage}[b]{\linewidth}\raggedright
Change
\end{minipage} &
\begin{minipage}[b]{\linewidth}\raggedright
Held-out
\end{minipage} &
\begin{minipage}[b]{\linewidth}\raggedright
Test
\end{minipage} &
\begin{minipage}[b]{\linewidth}\raggedright
Codabench
\end{minipage} &
\begin{minipage}[b]{\linewidth}\raggedright
Conclusion
\end{minipage} \\
\midrule\noalign{}
\endhead
\bottomrule\noalign{}
\endlastfoot
NeuralBench EEGNet, seed 34 & Official
baseline as a floor & braindecode
EEGNet, NeuralBench recipe & -- & 0.288
& \textbf{0.29} & Starting point
(submission 1001528) \\
EEGNet ensemble 34+35 & Seed averaging
helps & mean of 2
models\textquotesingle{} normalised
predictions & -- & 0.2975 & 0.30 & +1; 3
seeds tied (0.2970) \\
REVE frozen + ridge & A foundation
model\textquotesingle s features suffice
& official ridge probe on frozen REVE &
-- & 0.025 & -- & Chance level \\
REVE frozen + linear CLIP head & Same,
with a contrastive head & linear head,
CLIP loss & -- & 0.127 (120 Hz) / 0.086
(200 Hz) & -- & Far below EEGNet;
deprioritised \\
Our encoder (unseeded) & Bigger model +
higher lr + CLIP loss & §2.2 table, τ
0.07, 20 epochs & 0.4215 & 0.5415 &
\textbf{0.54} & +24 over baseline \\
3-encoder ensemble & Seeds help & 3
seeds & 0.4315 & 0.5530 & 0.55 & +1 \\
Control (seed 1) & Reference for round 1
& same, seed 1 & 0.4135 & 0.5450 & --
& \\
R1 dedup & False negatives hurt &
\texttt{dedup=True} & 0.4120 & -- & -- &
No gain (few duplicates per batch) \\
\textbf{R2 learnable τ} & τ = 0.07 is
too soft &
\texttt{learn\_temperature=True} &
\textbf{0.5175} & 0.6525 & \textbf{0.65}
& \textbf{+10.4}: keep \\
R3 τ + dedup + symmetric &
CLIP\textquotesingle s symmetric loss
helps & + \texttt{dedup},
\texttt{symmetric} & 0.4985 & -- & -- &
Worse than R2: drop \\
R4 R3 + augmentation & Invariance to
shifts/channels & ±50 ms circular shift,
10\% channel dropout & 0.3680 & -- & --
& Hurts a lot: drop \\
\textbf{R5 R2 + val data} & More data
helps & \texttt{-\/-use-val} (+25\%) &
\textbf{0.5525} & 0.7110 & \textbf{0.71}
& \textbf{+3.5}: keep \\
\#23 recalibration & Post-hoc linear map
fixes predictions & centering / ridge
maps on R2\textquotesingle s outputs &
0.376 / 0.046--0.052 (none: 0.5175) & --
& -- & All hurt: closed \\
\textbf{Final 3-seed ensemble} & All
concepts + ensemble & R5 recipe on all
668,400 windows, seeds 1--3 & n/a &
\textbf{0.7380} (single seeds
0.715--0.721) & \textbf{0.74} (top-1
0.40) & Current best (submission
1004941) \\
\textbf{R6 40 epochs} & 20 epochs were
too few & \texttt{n\_epochs=40}, early
stopping & \textbf{0.590} (epoch 36) &
not run & -- & \textbf{+3.75}: keep;
still rising at 40 \\
R15 floor 0.005 & τ is stuck at the 0.01
floor & \texttt{min\_temperature=0.005}
& 0.564 at epoch 31 (stopped at 32/40 by
a restart) & -- & -- & Trailed R6
throughout: no gain \\
R7 avg 1--4 reps (random) & Cleaner
training examples &
\texttt{-\/-avg-reps\ 4\ -\/-avg-reps-random}
& 0.555 (epoch 46, stopped at 51) & -- &
-- & Cut short by patience (§3.4);
faster per minute than R6; re-test with
scaled patience \\
R8 avg 2 reps & Same, fixed k = 2 &
\texttt{-\/-avg-reps\ 2} & 0.5535 (epoch
47, stopped at 52) & -- & -- & Same
caveat as R7 \\
R9 cosine & A decaying lr finds a better
optimum &
\texttt{lr\_schedule=\textquotesingle{}cosine\textquotesingle{}\ warmup\_epochs=1},
30 epochs, patience 8 & 0.570 (epoch 27)
& -- & -- & Below R6 (0.590 at 40
constant-lr epochs) \\
R10 batch 512 & More negatives per step
& R9 +
\texttt{-\/-batch-size\ 512\ lr=2e-3} &
0.581 (epoch 30) & -- & -- & ≈ R6 \\
R11 batch 1024 & Even more negatives &
R9 +
\texttt{-\/-batch-size\ 1024\ lr=4e-3} &
0.586 (epoch 27) & -- & -- & ≈ R6 in
\textbf{62 instead of 105 min}:
candidate with more epochs \\
\textbf{R12 wider convs} & The encoder
is capacity-limited &
\texttt{temporal\_filters=64\ spatial\_filters=128\ sep\_filters=256}
(2.26M params) & \textbf{0.615} (epoch
22, stopped at 27) & -- & -- &
\textbf{+2.5}: keep (5--6 points ahead
of R6 at every epoch);
\textasciitilde2.6× slower per epoch \\
R13 wider head & Same, in the head &
\texttt{hidden\_dim=1024\ head\_blocks=1}
& 0.529 at epoch 17 (R6 there: 0.545) &
-- & -- & Trailing: cancelled, as was
R14 (both) \\
\textbf{Final R6 recipe} & 40 epochs on
all concepts & R6 recipe, all concepts +
val, seeds 1--3, ensembled
(\texttt{final\_r6\_ens3}) & n/a &
\textbf{0.7760} (seeds 0.755--0.766) &
\textbf{0.78} (top-1 0.44) & Current
best (submission 1007631, rank 60);
trained at \texttt{95e45ea} \\
\end{longtable}
}

Top-1 subject-aggregated for the
uploaded models: 0.09 (EEGNet) → 0.22 →
0.34 (R2) → 0.37 (R5) → 0.40 (final
ensemble) → \textbf{0.44} (final R6
recipe).

Overnight Oct 21→22: a 3-seed final of
the R12 recipe (wider convolutions, 24
epochs, \texttt{logs/round6.sh}),
packaged only if it beats 0.776 on the
local test. Unseen-participant
experiments are in §10.

\begin{center}\rule{0.5\linewidth}{0.5pt}\end{center}

\subsection{6. Submission and inference
path}\label{submission-and-inference-path}

\begin{figure}[H]\centering
\begin{adjustbox}{max width=\linewidth, max totalheight=0.8\textheight}
\begin{tikzpicture}[>={Stealth[length=2mm]}, box/.style={draw, rounded corners=2pt, align=center, font=\footnotesize, inner sep=4pt, fill=black!4}, every edge quotes/.style={font=\scriptsize, fill=white, inner sep=1pt, align=center, auto}]
\graph [layered layout, grow=down, level sep=9mm, sibling sep=6mm, nodes={box}, edges={semithick}] {
  T [as={train\_solver.py \\ load\_model (fresh init) → fit → save\_model}];
  F [as={submissions/NAME/ \\ submission.py + weights.pt \\ (+ run\_config*.json)}];
  Z [as={make zip-submission SUBMISSION=submissions/NAME \\ (run\_config* left out)}];
  CB [as={Codabench: unzip read-only, \\ Solver.load\_model(meta), then predict(X) per batch}];
  SC [as={objective.evaluate\_result → leaderboard}];
  T ->[] F;
  F ->[] Z;
  Z ->[] CB;
  CB ->[] SC;
};
\end{tikzpicture}
\end{adjustbox}
\end{figure}

\begin{itemize}
\tightlist
\item
  \textbf{Contract}
  (\texttt{benchmark\_utils.base\_solver.CompetSolver}):
  Codabench never trains. It calls
  \texttt{Solver.load\_model(meta)} with
  \texttt{n\_chans}, \texttt{n\_times},
  \texttt{n\_outputs},
  \texttt{ch\_names},
  \texttt{chs\_info}, \texttt{device},
  \texttt{submission\_dir}, and then
  \texttt{model.predict(X)} on
  \texttt{(B,\ C,\ T)} batches;
  \texttt{predict} must return
  \texttt{(B,\ n\_outputs)}.
\item
  \textbf{Inference vs training in
  \texttt{load\_model}:} if
  \texttt{train\_loader} is set (a local
  training run) it builds a fresh,
  seeded model; otherwise it loads the
  shipped weights.
\item
  \textbf{Sizes from the weights
  (\texttt{encoder\_sizes}, \#27):} the
  number of filters, the head width and
  the number of residual blocks are read
  from the tensor shapes in the weights
  file, so the uploaded
  \texttt{submission.py} never needs
  editing after a wider run.
\item
  \textbf{Ensembles
  (\texttt{EmbeddingEnsemble}):} if the
  folder contains \texttt{weights\_*.pt}
  files, one model is built per file and
  their (unit-length) predictions are
  averaged.
\item
  \textbf{Run time on Codabench:}
  226--318 s per upload so far (limit 60
  min), e.g. 305 s for the 3-model
  ensemble.
\item
  \textbf{Reproducibility:} each run
  writes \texttt{run\_config.json}
  (command, commit, settings, seed,
  windows, minutes, best epoch) and
  \texttt{run\_config\_curve.json}
  (per-epoch loss, temperature, lr,
  held-out), which the top-3 audit
  needs.
\end{itemize}

\begin{center}\rule{0.5\linewidth}{0.5pt}\end{center}

\subsection{7. Paper map and
bibliography}\label{paper-map-and-bibliography}

\subsubsection{7.1 Idea → paper → where
in our
code}\label{idea-paper-where-in-our-code}

{\def\LTcaptype{none} % do not increment counter
\begin{longtable}[]{@{}
  >{\raggedright\arraybackslash}p{(\linewidth - 6\tabcolsep) * \real{0.2825}}
  >{\raggedright\arraybackslash}p{(\linewidth - 6\tabcolsep) * \real{0.2768}}
  >{\raggedright\arraybackslash}p{(\linewidth - 6\tabcolsep) * \real{0.2034}}
  >{\raggedright\arraybackslash}p{(\linewidth - 6\tabcolsep) * \real{0.2373}}@{}}
\toprule\noalign{}
\begin{minipage}[b]{\linewidth}\raggedright
Idea in our code
\end{minipage} &
\begin{minipage}[b]{\linewidth}\raggedright
Paper
\end{minipage} &
\begin{minipage}[b]{\linewidth}\raggedright
What to read
\end{minipage} &
\begin{minipage}[b]{\linewidth}\raggedright
Where it lives
\end{minipage} \\
\midrule\noalign{}
\endhead
\bottomrule\noalign{}
\endlastfoot
Temporal-then-spatial conv factorisation
& Schirrmeister et al. 2017 &
ShallowConvNet / DeepConvNet
architecture figures and their FBCSP
motivation &
\texttt{EEGToImageEncoder.temporal\_conv},
\texttt{.spatial\_conv} \\
Compact EEG CNN: block order, ELU,
pooling, dropout, separable conv &
Lawhern et al. 2018 (EEGNet) &
Architecture section and table of layers
& \texttt{EEGToImageEncoder} (whole conv
part) \\
Depthwise-separable convolution &
Chollet 2017 (Xception) & Section on
depthwise-separable convs &
\texttt{EEGToImageEncoder.sep\_conv} \\
BatchNorm & Ioffe \& Szegedy 2015 & &
every \texttt{BatchNorm*} \\
ELU / GELU activations & Clevert et al.
2016 / Hendrycks \& Gimpel 2016 & &
\texttt{spatial\_conv},
\texttt{sep\_conv} /
\texttt{projection\_head} \\
Dropout & Srivastava et al. 2014 & &
\texttt{Dropout(dropout)} \\
Contrastive loss with in-batch negatives
& van den Oord et al. 2018 (InfoNCE);
Radford et al. 2021 (CLIP) & CLIP §2.3
and its Figure 3 pseudo-code &
\texttt{clip\_loss} \\
Learnable temperature, logit scale ≤ 100
& Radford et al. 2021 (CLIP) & CLIP §2.5
(training details) & \texttt{Solver.fit}
(\texttt{log\_scale},
\texttt{min\_temperature}) \\
Symmetric loss & Radford et al. 2021 &
Figure 3 &
\texttt{clip\_loss(symmetric=True)} \\
CLIP loss for brain signals, retrieval
evaluation & Défossez et al. 2023;
Benchetrit et al. 2024 & Methods (loss,
retrieval metric) & \texttt{clip\_loss},
objective \\
Zero-shot EEG→image retrieval on
THINGS-EEG2 & Song et al. 2024 (NICE);
Li et al. 2024 (ATM) & Problem setup and
results tables & whole project \\
Residual projector + LayerNorm & Li et
al. 2024 (ATM); NICE\textquotesingle s
projector & Architecture section &
\texttt{ResidualBlock},
\texttt{head\_blocks} \\
DINOv2-giant target embeddings & Oquab
et al. 2023 (DINOv2) & Intro + feature
quality sections & targets (NeuralBench,
\texttt{facebook/dinov2-giant}) \\
AdamW & Loshchilov \& Hutter 2019 & &
\texttt{Solver.fit} optimiser \\
Cosine schedule & Loshchilov \& Hutter
2017 (SGDR) & &
\texttt{Solver.\_lr\_factor} \\
Warm-up, lr scaling with batch size &
Goyal et al. 2017 & §2.1--2.2 &
\texttt{\_lr\_factor}, R10/R11 \\
Repetition averaging / data protocol &
Gifford et al. 2022 (THINGS-EEG2) &
Paradigm and preprocessing &
\texttt{src/data/repetitions.py} \\
Concepts and images & Hebart et al. 2019
(THINGS) & &
\texttt{configs/heldout\_concepts.json} \\
Blurred-image targets (issue \#26) & UBP
(CVPR 2025) & & not implemented yet \\
Position-based spatial attention,
spatial dropout & Défossez et al. 2023 &
"Spatial attention" layer in Methods
(Fourier embedding of sensor positions)
& \texttt{SpatialAttention},
\texttt{Solver.virtual\_channels}
(§9.2) \\
Electrode positions (10-05 system) &
Oostenveld \& Praamstra 2001
\textbf{(unverified)};
MNE\textquotesingle s
\texttt{standard\_1005} montage & &
\texttt{ELECTRODE\_XY},
\texttt{channel\_positions} \\
Held-out participants (cross-subject
validation) & standard practice;
motivated by the sealed
final\textquotesingle s unseen subjects
& & \texttt{heldout\_subjects\_split}
(§9.1) \\
EEG foundation model: frozen +
channel-averaged (weak), then tokens +
attention-pooling head & El Ouahidi et
al. 2025 (REVE) & &
\texttt{submission/reve\_submission.py},
\texttt{scripts/cache\_reve\_tokens.py}
(§10.4) \\
More pretraining data matters most out
of distribution (motivates more
participants, REVE) & Kolesnikov et al.
2020 (BiT) \textbf{(unverified)} &
ObjectNet figure & §10.1 participant
curve \\
Participant-adversarial training (DANN,
domain confusion) & Ganin \& Lempitsky
2015; Tzeng et al. 2015; Elazar \&
Goldberg 2018 \textbf{(unverified)} & &
\texttt{SubjectAdversary},
\texttt{adversarial\_weight} (§10.3) \\
Final-phase data family & Xu et al. 2024
(Alljoined1); Xu et al. 2025
(Alljoined-1.6M) & Dataset papers &
issues \#4, \#21 \\
A convolution + transformer EEG encoder
(candidate for \#12) & Song et al. 2023
(EEG Conformer) & Architecture figure &
not implemented \\
\end{longtable}
}

\subsubsection{7.2
Bibliography}\label{bibliography}

Verification: "verified" means the
citation was checked against metadata
shipped in our installed packages
(braindecode, NeuralFetch, PyTorch,
transformers). Web search was not
available while writing this, so the
others are from memory and marked
\textbf{(unverified)}.

\textbf{Datasets}

\begin{itemize}
\tightlist
\item
  Gifford, A. T., Dwivedi, K., Roig, G.,
  \& Cichy, R. M. (2022). A large and
  rich EEG dataset for modeling human
  visual object recognition.
  \emph{NeuroImage}.
  \url{doi:10.1016/j.neuroimage.2022.119754}
  --- verified (NeuralFetch
  \texttt{gifford2022large}). Data:
  osf.io/3jk45.
\item
  Hebart, M. N., et al. (2019). THINGS:
  A database of 1,854 object concepts
  and more than 26,000 naturalistic
  object images. \emph{PLoS ONE} 14(10):
  e0223792. \textbf{(unverified)}
\item
  Xu, J., Aristimunha, B., Feucht, M.
  E., Qian, E., et al. (2024). Alljoined
  --- A dataset for EEG-to-Image
  decoding. arXiv:2404.05553 ---
  verified (NeuralFetch
  \texttt{xu2024alljoined}; 8 subjects).
\item
  Xu, J., Nunes, U. B., Jiang, W.,
  Ryther, S., Pringle, J., Scotti, P.
  S., Delorme, A., \& Kneeland, R.
  (2025). Alljoined-1.6M: A
  million-trial EEG-image dataset for
  evaluating affordable brain-computer
  interfaces. arXiv:2508.18571 ---
  verified (NeuralFetch
  \texttt{xu2025alljoined}; 20 subjects,
  Emotiv headset).
\end{itemize}

\textbf{EEG architectures}

\begin{itemize}
\tightlist
\item
  Lawhern, V. J., Solon, A. J.,
  Waytowich, N. R., Gordon, S. M., Hung,
  C. P., \& Lance, B. J. (2018). EEGNet:
  a compact convolutional neural network
  for EEG-based brain--computer
  interfaces. \emph{Journal of Neural
  Engineering} 15(5), 056013.
  arXiv:1611.08024 --- verified
  (braindecode).
\item
  Schirrmeister, R. T., Springenberg, J.
  T., Fiederer, L. D. J., Glasstetter,
  M., Eggensperger, K., Tangermann, M.,
  Hutter, F., \& Ball, T. (2017). Deep
  learning with convolutional neural
  networks for EEG decoding and
  visualization. \emph{Human Brain
  Mapping}. \url{doi:10.1002/hbm.23730}
  --- verified (braindecode).
\item
  Song, Y., Zheng, Q., Liu, B., \& Gao,
  X. (2023). EEG Conformer:
  Convolutional transformer for EEG
  decoding and visualization. \emph{IEEE
  Transactions on Neural Systems and
  Rehabilitation Engineering} 31,
  710--719 --- verified (braindecode;
  braindecode dates it 2022, the journal
  volume is 2023).
\item
  Chollet, F. (2017). Xception: Deep
  learning with depthwise separable
  convolutions. \emph{CVPR} --- verified
  (cited in
  braindecode\textquotesingle s EEGNet).
\item
  El Ouahidi, Y., Lys, J., Thölke, P.,
  Farrugia, N., Pasdeloup, B., Gripon,
  V., Jerbi, K., \& Lioi, G. (2025).
  REVE: A foundation model for EEG ---
  adapting to any setup with large-scale
  pretraining on 25,000 subjects.
  \emph{NeurIPS 2025}.
  openreview.net/forum?id=ZeFMtRBy4Z ---
  verified (braindecode).
\end{itemize}

\textbf{EEG/MEG → image decoding}

\begin{itemize}
\tightlist
\item
  Song, Y., Liu, B., Li, X., Shi, N.,
  Wang, Y., \& Gao, X. (2024). Decoding
  natural images from EEG for object
  recognition (NICE). \emph{ICLR 2024}.
  arXiv:2308.13234 \textbf{(unverified)}
\item
  Li, D., Wei, C., Li, S., Zou, J., \&
  Liu, Q. (2024). Visual decoding and
  reconstruction via EEG embeddings with
  guided diffusion (ATM). \emph{NeurIPS
  2024}. arXiv:2403.07721
  \textbf{(unverified)}
\item
  "Bridging the vision-brain gap with an
  uncertainty-aware blur prior" (UBP).
  \emph{CVPR 2025}. \textbf{(unverified:
  authors and arXiv ID; search the
  title)}
\item
  Benchetrit, Y., Banville, H., \& King,
  J.-R. (2024). Brain decoding: toward
  real-time reconstruction of visual
  perception. \emph{ICLR 2024}.
  arXiv:2310.19812 \textbf{(unverified)}
\item
  Défossez, A., Caucheteux, C., Rapin,
  J., Kabeli, O., \& King, J.-R. (2023).
  Decoding speech perception from
  non-invasive brain recordings.
  \emph{Nature Machine Intelligence}
  5(10), 1097--1107 --- verified
  (braindecode REVE references).
  arXiv:2208.12266 \textbf{(ID
  unverified)}
\end{itemize}

\textbf{Representation learning and
optimisation}

\begin{itemize}
\tightlist
\item
  Radford, A., et al. (2021). Learning
  transferable visual models from
  natural language supervision (CLIP).
  \emph{ICML 2021}. arXiv:2103.00020
  \textbf{(unverified)}
\item
  van den Oord, A., Li, Y., \& Vinyals,
  O. (2018). Representation learning
  with contrastive predictive coding.
  arXiv:1807.03748 \textbf{(unverified)}
\item
  Oquab, M., et al. (2023). DINOv2:
  Learning robust visual features
  without supervision. arXiv:2304.07193
  --- verified
  (transformers\textquotesingle{} DINOv2
  module).
\item
  Loshchilov, I., \& Hutter, F. (2019).
  Decoupled weight decay regularization
  (AdamW). \emph{ICLR 2019}.
  arXiv:1711.05101 --- verified (PyTorch
  \texttt{AdamW} docstring).
\item
  Loshchilov, I., \& Hutter, F. (2017).
  SGDR: Stochastic gradient descent with
  warm restarts. \emph{ICLR 2017}.
  arXiv:1608.03983 --- verified (PyTorch
  \texttt{lr\_scheduler}).
\item
  Goyal, P., et al. (2017). Accurate,
  large minibatch SGD: Training ImageNet
  in 1 hour. arXiv:1706.02677
  \textbf{(unverified)}
\item
  Ioffe, S., \& Szegedy, C. (2015).
  Batch normalization. \emph{ICML 2015}.
  \textbf{(unverified)}
\item
  Clevert, D.-A., Unterthiner, T., \&
  Hochreiter, S. (2016). Fast and
  accurate deep network learning by
  exponential linear units (ELUs).
  \emph{ICLR 2016}.
  \textbf{(unverified)}
\item
  Hendrycks, D., \& Gimpel, K. (2016).
  Gaussian error linear units (GELUs).
  arXiv. \textbf{(unverified)}
\item
  Srivastava, N., et al. (2014).
  Dropout: A simple way to prevent
  neural networks from overfitting.
  \emph{JMLR}. \textbf{(unverified)}
\end{itemize}

\textbf{Added Oct 21 (from memory,
unverified: check titles before citing)}

\begin{itemize}
\tightlist
\item
  Kolesnikov, A., Beyer, L., Zhai, X.,
  Puigcerver, J., Yung, J., Gelly, S.,
  \& Houlsby, N. (2020). Big Transfer
  (BiT): General Visual Representation
  Learning. ECCV. \textbf{(unverified)}
\item
  Ganin, Y., \& Lempitsky, V. (2015).
  Unsupervised Domain Adaptation by
  Backpropagation. ICML (the
  gradient-reversal layer, DANN).
  \textbf{(unverified)}
\item
  Tzeng, E., Hoffman, J., Darrell, T.,
  \& Saenko, K. (2015). Simultaneous
  Deep Transfer Across Domains and
  Tasks. ICCV (domain-confusion loss).
  \textbf{(unverified)}
\item
  Elazar, Y., \& Goldberg, Y. (2018).
  Adversarial Removal of Demographic
  Attributes from Text Data. EMNLP
  (adversarially "removed" attributes
  stay recoverable).
  \textbf{(unverified)}
\end{itemize}

\begin{center}\rule{0.5\linewidth}{0.5pt}\end{center}

\subsection{8. Teaching
kit}\label{teaching-kit}

\subsubsection{8.1 A 60--90 minute
session}\label{a-6090-minute-session}

{\def\LTcaptype{none} % do not increment counter
\begin{longtable}[]{@{}
  >{\raggedright\arraybackslash}p{(\linewidth - 4\tabcolsep) * \real{0.0706}}
  >{\raggedright\arraybackslash}p{(\linewidth - 4\tabcolsep) * \real{0.5882}}
  >{\raggedright\arraybackslash}p{(\linewidth - 4\tabcolsep) * \real{0.3412}}@{}}
\toprule\noalign{}
\begin{minipage}[b]{\linewidth}\raggedright
Time
\end{minipage} &
\begin{minipage}[b]{\linewidth}\raggedright
Topic
\end{minipage} &
\begin{minipage}[b]{\linewidth}\raggedright
Show
\end{minipage} \\
\midrule\noalign{}
\endhead
\bottomrule\noalign{}
\endlastfoot
0--10 min & The task and why
it\textquotesingle s retrieval, not
classification; chance level 2.5\% &
how\_it\_works.md, first diagram \\
10--20 & The data: THINGS-EEG2, windows,
why repetitions matter (4 vs 80) &
how\_it\_works.md §2; §4.2 here \\
20--35 & The encoder, layer by layer:
temporal → spatial → separable → head;
receptive field & §1.1 diagram and §1.2
table \\
35--50 & The loss: the B×B matrix,
negatives, temperature; why not
regression & §3.1 formula;
how\_it\_works.md loss diagram \\
50--60 & Evaluation and the held-out
split; why held-out \textless{} test &
§4.1 diagram, §4.3 table \\
60--75 & The experiment log: what worked
(τ, data, epochs), what
didn\textquotesingle t, and why & §5
table \\
75--90 & What\textquotesingle s next:
final phase = 32 channels + unseen
people; headset-independent spatial
layer; held-out participants &
how\_it\_works.md §8 \\
\end{longtable}
}

Tip: before the session, run exercise 2
so the scores appear live on screen.

\subsubsection{8.2
Check-your-understanding questions (with
answers)}\label{check-your-understanding-questions-with-answers}

\begin{enumerate}
\def\labelenumi{\arabic{enumi}.}
\tightlist
\item
  \textbf{Why can the model recognise
  images of concepts it never trained
  on?} It predicts a point in
  DINOv2\textquotesingle s embedding
  space, not a class. DINOv2 places any
  image, including unseen ones, in that
  space, so retrieval works for new
  images (zero-shot).
\item
  \textbf{What is chance-level top-5
  accuracy on the test set?} 5/200 =
  2.5\%.
\item
  \textbf{Why is subject-aggregated
  accuracy (0.74) so much higher than
  pooled (0.31)?} It averages
  \textasciitilde80 predictions per
  (subject, image) before ranking, which
  cancels most of the single-trial
  noise.
\item
  \textbf{What does the temporal
  convolution learn, in
  signal-processing terms?} 32 FIR
  filters of 258 ms shared across
  electrodes: band-pass filters or
  response-shape detectors.
\item
  \textbf{Why does the model have to be
  retrained for the 32-channel final
  phase?} The spatial
  conv\textquotesingle s kernel spans
  exactly 63 electrodes in a fixed
  order; its weights have no meaning for
  a different montage.
\item
  \textbf{Where are 90\% of the
  parameters?} In the projection head
  (896 → 512 → 1536).
\item
  \textbf{Write the InfoNCE loss for one
  row and say what the negatives are.}
  \(-\log \frac{e^{s_{ii}/\tau}}{\sum_j e^{s_{ij}/\tau}}\);
  the negatives are the other images in
  the same batch.
\item
  \textbf{What does lowering the
  temperature do?} Sharpens the softmax
  so the loss focuses on the hardest
  wrong images. Learning it gave +10
  held-out; going below 0.01 gave
  nothing.
\item
  \textbf{Why did \texttt{dedup} make no
  difference?} Same-image duplicates are
  rare in a batch of 256 (about 2
  pairs), so there were few false
  negatives to remove.
\item
  \textbf{Give one plausible reason
  augmentation by time shift hurt.} The
  flattened readout relies on absolute
  timing of responses; shifting by up to
  ±50 ms (with wrap-around) corrupts
  that.
\item
  \textbf{Why is the held-out score
  lower than the test score for the same
  model?} Held-out images were seen
  \textasciitilde4 times per subject,
  test images \textasciitilde80, so
  subject aggregation removes less noise
  on held-out.
\item
  \textbf{Why don\textquotesingle t we
  choose models by their test score?}
  The test set is the public leaderboard
  data; selecting on it overfits to it
  and tells us nothing about the sealed
  phase. We select on the held-out
  split.
\item
  \textbf{Why do we retrain on all
  concepts for uploads?} More concepts
  cover DINOv2\textquotesingle s space
  better; it added \textasciitilde+2.7
  on test together with the 3-seed
  ensemble.
\item
  \textbf{Why was R7 (repetition
  averaging) not a fair loss?} Its
  epochs are 4× shorter, so patience 5
  stopped it while it was still
  improving.
\item
  \textbf{What does
  \texttt{encoder\_sizes} solve?} It
  reads the layer sizes from the
  weights, so the same
  \texttt{submission.py} can load a
  model of any width without edits.
\end{enumerate}

\subsubsection{8.3 Hands-on
exercises}\label{hands-on-exercises}

\textbf{Exercise 1: count the parameters
and check the shapes (CPU, 1 minute).}

\begin{Shaded}
\begin{Highlighting}[]
\CommentTok{\# save as count\_params.py at the repo root, run: .venv/bin/python count\_params.py}
\ImportTok{import}\NormalTok{ importlib.util, sys, torch}
\NormalTok{sys.path.insert(}\DecValTok{0}\NormalTok{, }\StringTok{"external/2026{-}competition/tracks/image\_decoding"}\NormalTok{)}
\NormalTok{spec }\OperatorTok{=}\NormalTok{ importlib.util.spec\_from\_file\_location(}\StringTok{"s"}\NormalTok{, }\StringTok{"submission/submission.py"}\NormalTok{)}
\NormalTok{m }\OperatorTok{=}\NormalTok{ importlib.util.module\_from\_spec(spec)}\OperatorTok{;}\NormalTok{ spec.loader.exec\_module(m)}
\NormalTok{enc }\OperatorTok{=}\NormalTok{ m.EEGToImageEncoder(n\_chans}\OperatorTok{=}\DecValTok{63}\NormalTok{, n\_times}\OperatorTok{=}\DecValTok{120}\NormalTok{, output\_dim}\OperatorTok{=}\DecValTok{1536}\NormalTok{).}\BuiltInTok{eval}\NormalTok{()}
\NormalTok{x }\OperatorTok{=}\NormalTok{ torch.zeros(}\DecValTok{2}\NormalTok{, }\DecValTok{1}\NormalTok{, }\DecValTok{63}\NormalTok{, }\DecValTok{120}\NormalTok{)}
\ControlFlowTok{for}\NormalTok{ block }\KeywordTok{in}\NormalTok{ (enc.temporal\_conv, enc.spatial\_conv, enc.sep\_conv):}
    \ControlFlowTok{for}\NormalTok{ layer }\KeywordTok{in}\NormalTok{ block:}
\NormalTok{        y }\OperatorTok{=}\NormalTok{ layer(x)}
        \BuiltInTok{print}\NormalTok{(}\SpecialStringTok{f"}\SpecialCharTok{\{}\BuiltInTok{str}\NormalTok{(layer)[:}\DecValTok{60}\NormalTok{]}\SpecialCharTok{:60s\}}\SpecialStringTok{ }\SpecialCharTok{\{}\BuiltInTok{tuple}\NormalTok{(x.shape)}\SpecialCharTok{\}}\SpecialStringTok{ {-}\textgreater{} }\SpecialCharTok{\{}\BuiltInTok{tuple}\NormalTok{(y.shape)}\SpecialCharTok{\}}\SpecialStringTok{ "}
              \SpecialStringTok{f"}\SpecialCharTok{\{}\BuiltInTok{sum}\NormalTok{(p.numel() }\ControlFlowTok{for}\NormalTok{ p }\KeywordTok{in}\NormalTok{ layer.parameters())}\SpecialCharTok{\}}\SpecialStringTok{"}\NormalTok{)}
\NormalTok{        x }\OperatorTok{=}\NormalTok{ y}
\BuiltInTok{print}\NormalTok{(}\StringTok{"total:"}\NormalTok{, }\BuiltInTok{sum}\NormalTok{(p.numel() }\ControlFlowTok{for}\NormalTok{ p }\KeywordTok{in}\NormalTok{ enc.parameters()))  }\CommentTok{\# 1,387,872}
\end{Highlighting}
\end{Shaded}

Questions: where do 129,024 parameters
come from? (32 × 63 × 64.) What changes
with
\texttt{EEGToImageEncoder(32,\ 120,\ 1536)},
the final phase\textquotesingle s
channel count?

\textbf{Exercise 2: score an existing
model like Codabench does (≈1 min with
the cache).}

\begin{Shaded}
\begin{Highlighting}[]
\ExtensionTok{python}\NormalTok{ scripts/eval\_local.py submissions/e6\_r5\_learntemp\_useval }\AttributeTok{{-}{-}cache}\NormalTok{ cache/windows/things\_eeg2}
\CommentTok{\# expect 0.7110 / 0.3745 / 0.2824 / 0.0958 (R5, Codabench 0.71)}
\ExtensionTok{python}\NormalTok{ scripts/eval\_local.py submissions/e6\_r5\_learntemp\_useval }\AttributeTok{{-}{-}split}\NormalTok{ heldout }\AttributeTok{{-}{-}cache}\NormalTok{ cache/windows/things\_eeg2}
\CommentTok{\# expect 0.5525 held{-}out top{-}5 subject{-}aggregated}
\end{Highlighting}
\end{Shaded}

Questions: why is the second number
lower? Which column is the leaderboard
metric?

\textbf{Exercise 3: plot learning curves
(CPU, 2 minutes).}

\begin{Shaded}
\begin{Highlighting}[]
\CommentTok{\# save as plot\_curves.py, run: .venv/bin/python plot\_curves.py}
\ImportTok{import}\NormalTok{ json, matplotlib.pyplot }\ImportTok{as}\NormalTok{ plt}
\NormalTok{fig, (a, b) }\OperatorTok{=}\NormalTok{ plt.subplots(}\DecValTok{1}\NormalTok{, }\DecValTok{2}\NormalTok{, figsize}\OperatorTok{=}\NormalTok{(}\DecValTok{10}\NormalTok{, }\DecValTok{4}\NormalTok{))}
\ControlFlowTok{for}\NormalTok{ run }\KeywordTok{in}\NormalTok{ [}\StringTok{"m2\_r6\_epochs40"}\NormalTok{, }\StringTok{"m2\_r7\_avgreps4rand"}\NormalTok{]:}
\NormalTok{    curve }\OperatorTok{=}\NormalTok{ json.load(}\BuiltInTok{open}\NormalTok{(}\SpecialStringTok{f"submissions/}\SpecialCharTok{\{}\NormalTok{run}\SpecialCharTok{\}}\SpecialStringTok{/run\_config\_curve.json"}\NormalTok{))}
\NormalTok{    ep }\OperatorTok{=}\NormalTok{ [r[}\StringTok{"epoch"}\NormalTok{] }\ControlFlowTok{for}\NormalTok{ r }\KeywordTok{in}\NormalTok{ curve]}
\NormalTok{    a.plot(ep, [r[}\StringTok{"heldout"}\NormalTok{] }\ControlFlowTok{for}\NormalTok{ r }\KeywordTok{in}\NormalTok{ curve], label}\OperatorTok{=}\NormalTok{run)}
\NormalTok{    b.plot(ep, [r[}\StringTok{"temperature"}\NormalTok{] }\ControlFlowTok{for}\NormalTok{ r }\KeywordTok{in}\NormalTok{ curve], label}\OperatorTok{=}\NormalTok{run)}
\NormalTok{a.}\BuiltInTok{set}\NormalTok{(xlabel}\OperatorTok{=}\StringTok{"epoch"}\NormalTok{, ylabel}\OperatorTok{=}\StringTok{"held{-}out top{-}5 (subject{-}agg)"}\NormalTok{)}\OperatorTok{;}\NormalTok{ b.}\BuiltInTok{set}\NormalTok{(xlabel}\OperatorTok{=}\StringTok{"epoch"}\NormalTok{, ylabel}\OperatorTok{=}\StringTok{"temperature"}\NormalTok{, yscale}\OperatorTok{=}\StringTok{"log"}\NormalTok{)}
\NormalTok{a.legend()}\OperatorTok{;}\NormalTok{ plt.tight\_layout()}\OperatorTok{;}\NormalTok{ plt.savefig(}\StringTok{"curves.png"}\NormalTok{)}\OperatorTok{;} \BuiltInTok{print}\NormalTok{(}\StringTok{"wrote curves.png"}\NormalTok{)}
\end{Highlighting}
\end{Shaded}

Questions: when does the temperature
reach its 0.01 floor? Re-plot with the
x-axis in training samples (epochs ×
587,574 for R6, × 145,400 for R7): which
run learns faster per sample?

\begin{center}\rule{0.5\linewidth}{0.5pt}\end{center}

\subsection{9. Next: new people and new
headsets (in
progress)}\label{next-new-people-and-new-headsets-in-progress}

Everything above was tuned on
THINGS-EEG2, where the test participants
are the \textbf{same 10 people} the
model trained on, recorded with
\textbf{63 channels}. The sealed final
is different: \textbf{11 new people, a
32-channel Emotiv headset, new images.}
Two pieces of code (commits
\texttt{25b12fb} and \texttt{d3ed4ea},
Oct 20) prepare for that. Neither has
been trained at scale yet.

\subsubsection{9.1 Measuring what the
final measures: held-out
participants}\label{measuring-what-the-final-measures-held-out-participants}

Our held-out split (§4.2) holds out
\emph{concepts} but not \emph{people},
so it can\textquotesingle t tell us how
well a model works on someone it has
never seen, which is notoriously the
hard part of EEG (skull thickness,
cortical folding and cap placement
differ between people).

\texttt{train\_solver.py\ -\/-holdout-subjects}
/
\texttt{eval\_local.py\ -\/-split\ heldout\_subjects}
(\texttt{heldout\_subjects\_split} in
\texttt{src/data/heldout.py}):

\begin{itemize}
\tightlist
\item
  the 200 held-out concepts are removed
  from training, as before;
\item
  \textbf{participants 6 and 7}
  (\texttt{configs/heldout\_subjects.json},
  random with seed 0) are removed
  entirely: training uses 8 people;
\item
  the score is computed on those 2
  people\textquotesingle s windows of
  the held-out concepts. Each held-out
  concept has 10 images, so there are
  \textbf{10 galleries of 200 images}
  (gallery \emph{g} = the \emph{g}-th
  image of every concept), each scored
  with the official metric, then
  averaged.
\end{itemize}

Why 10 galleries: one gallery with 2
people gives only 400 subject-aggregated
queries of \textasciitilde4 repetitions
each (standard error ≈ 2.5 points, too
noisy to judge 1--2-point changes). Ten
galleries give 16,174 windows,
\textasciitilde4,000 queries, and a
measured standard error of \textbf{≈ 0.8
points} (\texttt{eval\_local} prints it
as
\texttt{top5\_acc\_subject\_agg\_se}).

{\def\LTcaptype{none} % do not increment counter
\begin{longtable}[]{@{}
  >{\raggedright\arraybackslash}p{(\linewidth - 4\tabcolsep) * \real{0.2738}}
  >{\raggedright\arraybackslash}p{(\linewidth - 4\tabcolsep) * \real{0.3571}}
  >{\raggedright\arraybackslash}p{(\linewidth - 4\tabcolsep) * \real{0.3690}}@{}}
\toprule\noalign{}
\begin{minipage}[b]{\linewidth}\raggedright
Model
\end{minipage} &
\begin{minipage}[b]{\linewidth}\raggedright
Trained on participants 6 and 7?
\end{minipage} &
\begin{minipage}[b]{\linewidth}\raggedright
Held-out-participants score (top-5 SA)
\end{minipage} \\
\midrule\noalign{}
\endhead
\bottomrule\noalign{}
\endlastfoot
R6 & yes (reference, not a real held-out
score) & 0.6352 ± 0.0077 \\
\textbf{R16} (R6 recipe,
\texttt{-\/-holdout-subjects}) & no &
\textbf{0.2333} (best epoch 16 of 21) \\
\end{longtable}
}

Unseen people cost about 40 points. R16
also peaked much earlier (epoch 16) than
R6 on seen people (epoch 36): longer
training fits the training participants,
not people in general.

\subsubsection{9.2 One model for any
headset: spatial
attention}\label{one-model-for-any-headset-spatial-attention}

Our \texttt{spatial\_conv} (§1.3) is a
convolution with kernel height = number
of channels, in a fixed order. Its
weights only make sense for
THINGS-EEG2\textquotesingle s 63
channels, so a THINGS-EEG2 model
can\textquotesingle t read
Emotiv\textquotesingle s 32 channels,
and the two datasets
can\textquotesingle t be trained
together without remapping.

Défossez et al. (2023) solved this for
MEG with a \textbf{spatial attention}
layer, which we adapt
(\texttt{SpatialAttention}, enabled with
\texttt{-\/-set\ virtual\_channels=64}):

\begin{figure}[H]\centering
\begin{adjustbox}{max width=\linewidth, max totalheight=0.8\textheight}
\begin{tikzpicture}[>={Stealth[length=2mm]}, box/.style={draw, rounded corners=2pt, align=center, font=\footnotesize, inner sep=4pt, fill=black!4}, every edge quotes/.style={font=\scriptsize, fill=white, inner sep=1pt, align=center, auto}]
\graph [layered layout, grow=down, level sep=9mm, sibling sep=6mm, nodes={box}, edges={semithick}] {
  N [as={Channel names \\ (meta['ch\_names'])}];
  P [as={2D scalp position of each electrode \\ (standard 10-05 table, ELECTRODE\_XY)}];
  F [as={Fourier features of the position \\ cos/sin(2π(kx·x + ky·y)), 16×16 frequencies \\ → 512 numbers per electrode}];
  S [as={Score of electrode c for virtual channel j \\ = learned weights\_j · features\_c}];
  SM [as={softmax over electrodes \\ (unknown names masked out; \\ training: drop electrodes near a random point)}];
  X [as={EEG (B, C, 120) \\ C = 63 or 32}];
  M [as={Each virtual channel = weighted \\ average of the electrodes}];
  V [as={(B, 64, 120): same shape for any headset}];
  E [as={the existing encoder \\ (temporal conv → spatial conv → …)}];
  N ->[] P;
  P ->[] F;
  F ->[] S;
  S ->[] SM;
  X ->[] M;
  SM ->[] M;
  M ->[] V;
  V ->[] E;
};
\end{tikzpicture}
\end{adjustbox}
\end{figure}

\begin{itemize}
\tightlist
\item
  \textbf{Why it\textquotesingle s
  montage-independent:} the learned
  parameters (\texttt{heads}, 64 × 512 ≈
  33k numbers) define a smooth
  \emph{function of scalp position}, not
  one weight per channel. A new headset
  just evaluates that function at its
  own electrode positions.
\item
  \textbf{Positions:} the pipeline only
  passes electrode \emph{names}, and
  Codabench\textquotesingle s worker
  only has \texttt{torch}, so
  \texttt{submission.py} embeds a table
  of 343 standard 10-05 positions,
  generated once from
  MNE\textquotesingle s
  \texttt{standard\_1005} montage
  (sphere-fitted, then flattened by an
  azimuthal-equidistant projection from
  the top of the head, like a
  topographic map), at a fixed scale, so
  63- and 32-channel headsets share one
  coordinate frame.
\item
  \textbf{Spatial dropout}
  (\texttt{spatial\_dropout=0.15}):
  during training, every electrode
  within that radius of a random scalp
  point is ignored, so the model learns
  not to depend on any one region, which
  is useful when the
  final\textquotesingle s headset lacks
  some of THINGS-EEG2\textquotesingle s
  electrodes.
\item
  \textbf{Checked} (synthetic data):
  with \texttt{virtual\_channels=0} (the
  default) training is bit-identical to
  before; a model trained on 63 channels
  reloads to identical predictions and
  runs on a 32-channel Emotiv montage
  with the same weights. Positions are
  per-montage buffers, not saved in
  \texttt{weights.pt}.
\end{itemize}

\textbf{Result (R17 = R16 +
\texttt{virtual\_channels=64\ spatial\_dropout=0.15}):}
best 0.2553 vs 0.2333, and +1.6 on
average over the last 10 epochs: a
small, consistent gain (one seed). But
R17 trained \textbf{3.2× slower} (116 vs
36 min). It isn\textquotesingle t the
softmax/einsum layout: with the layer in
front, the first convolution must also
compute the gradient of its input, which
looks slow on the Apple GPU for this
shape; a clean benchmark needs an idle
GPU (to do).

\subsubsection{9.3 The plan (as of Oct
20; see §10 for what
happened)}\label{the-plan-as-of-oct-20-see-10-for-what-happened}

\begin{enumerate}
\def\labelenumi{\arabic{enumi}.}
\tightlist
\item
  Train the R6 recipe with
  \texttt{-\/-holdout-subjects} (the
  unseen-people baseline).
\item
  Same with
  \texttt{-\/-set\ virtual\_channels=64\ spatial\_dropout=0.15}:
  does spatial attention cost anything
  on THINGS-EEG2, and does it help on
  unseen people?
\item
  Then joint training on THINGS-EEG2 +
  Alljoined (issue \#21), possible only
  with step 2\textquotesingle s layer,
  validated on
  Alljoined\textquotesingle s held-out
  participants (issue \#4).
\end{enumerate}

\begin{center}\rule{0.5\linewidth}{0.5pt}\end{center}

\subsection{10. Unseen participants and
Alljoined: results and new tools (Oct
21)}\label{unseen-participants-and-alljoined-results-and-new-tools-oct-21}

\subsubsection{10.1 What helps on unseen
people (THINGS-EEG2, round
3)}\label{what-helps-on-unseen-people-things-eeg2-round-3}

Each run is R16 (0.2333) with
\textbf{one} change, scored on
participants 6 and 7 (10 galleries,
standard error ≈ 0.8 points):

{\def\LTcaptype{none} % do not increment counter
\begin{longtable}[]{@{}
  >{\raggedright\arraybackslash}p{(\linewidth - 8\tabcolsep) * \real{0.0508}}
  >{\raggedright\arraybackslash}p{(\linewidth - 8\tabcolsep) * \real{0.4153}}
  >{\raggedright\arraybackslash}p{(\linewidth - 8\tabcolsep) * \real{0.0593}}
  >{\raggedright\arraybackslash}p{(\linewidth - 8\tabcolsep) * \real{0.0508}}
  >{\raggedright\arraybackslash}p{(\linewidth - 8\tabcolsep) * \real{0.4237}}@{}}
\toprule\noalign{}
\begin{minipage}[b]{\linewidth}\raggedright
Run
\end{minipage} &
\begin{minipage}[b]{\linewidth}\raggedright
Change
\end{minipage} &
\begin{minipage}[b]{\linewidth}\raggedright
Score
\end{minipage} &
\begin{minipage}[b]{\linewidth}\raggedright
Δ
\end{minipage} &
\begin{minipage}[b]{\linewidth}\raggedright
Verdict
\end{minipage} \\
\midrule\noalign{}
\endhead
\bottomrule\noalign{}
\endlastfoot
\textbf{R18} &
\texttt{window\_norm=True}: each channel
standardised over time inside each
window & \textbf{0.274} & \textbf{+4.1}
& ✓ keep: removes per-person
amplitude/offset differences, costs
nothing. REVE\textquotesingle s own
preprocessing does the same per-window
z-scoring \\
R20 & wider convolutions
(R12\textquotesingle s) & 0.253 & +2.0 &
✓ helps, smaller than on seen people
(+2.5), 3.5× slower \\
R21 &
\texttt{max\_shift=6\ channel\_dropout=0.1}
& 0.230 & −0.3 & -- no effect (it hurt
on seen people, R4) \\
R19 & \texttt{adversarial\_weight=1.0}
(domain confusion, §10.3) & 0.190 & −4.3
& ✗ hurts: drop \\
\end{longtable}
}

\textbf{The participant curve}
(\texttt{-\/-train-subjects\ N}: the
first N of the 8 training participants):

{\def\LTcaptype{none} % do not increment counter
\begin{longtable}[]{@{}
  >{\raggedright\arraybackslash}p{(\linewidth - 8\tabcolsep) * \real{0.4545}}
  >{\raggedright\arraybackslash}p{(\linewidth - 8\tabcolsep) * \real{0.1364}}
  >{\raggedright\arraybackslash}p{(\linewidth - 8\tabcolsep) * \real{0.1364}}
  >{\raggedright\arraybackslash}p{(\linewidth - 8\tabcolsep) * \real{0.1364}}
  >{\raggedright\arraybackslash}p{(\linewidth - 8\tabcolsep) * \real{0.1364}}@{}}
\toprule\noalign{}
\begin{minipage}[b]{\linewidth}\raggedright
Training participants
\end{minipage} &
\begin{minipage}[b]{\linewidth}\raggedright
2
\end{minipage} &
\begin{minipage}[b]{\linewidth}\raggedright
4
\end{minipage} &
\begin{minipage}[b]{\linewidth}\raggedright
6
\end{minipage} &
\begin{minipage}[b]{\linewidth}\raggedright
8
\end{minipage} \\
\midrule\noalign{}
\endhead
\bottomrule\noalign{}
\endlastfoot
Unseen-people score & 0.161 & 0.183 &
0.187* & 0.233 \\
\end{longtable}
}

* stopped at epoch 8 (best epoch 3):
unreliable. The trend, about +3.6 points
per doubling and not flattening, is the
EEG version of the "more pretraining
data" lesson of BiT (Kolesnikov et al.,
2020): \textbf{more people is the
biggest lever.} Which people matters too
(one seed per point).

\subsubsection{10.2 Alljoined-1.6M: the
final\textquotesingle s
format}\label{alljoined-1.6m-the-finals-format}

\begin{itemize}
\tightlist
\item
  \textbf{Data:} 20 participants,
  32-channel Emotiv (Cz, FCz, AFz,
  Fp1/2, F1/2/5/6, CP1--6, P1--8,
  PO3/4/7/8, O1/2, Oz, POz, Pz), 120 Hz,
  \textasciitilde1.66M windows (train
  1,065,150 · val 270,466 · test
  324,294). NeuralFetch study
  \texttt{Xu2025Alljoined} (licence
  CC-BY-NC-ND-4.0,
  \texttt{ALLJOINED\_ACCEPT\_LICENCE=1}).
\item
  \textbf{Same images as THINGS-EEG2:}
  16,540 training + 200 test images; the
  200 test images\textquotesingle{}
  DINOv2 embeddings match
  THINGS-EEG2\textquotesingle s
  one-to-one (cosine 1.000000). Files
  are numbered (\texttt{stim\_file}),
  without concept names. Consequences:
  joint training is aligned, and
  \textbf{Alljoined\textquotesingle s
  test images are the warm-up
  leaderboard\textquotesingle s}, so its
  test split is never trained on.
\item
  \textbf{Validation split}
  (\texttt{-\/-dataset\ xu2025alljoined\ -\/-holdout-subjects}):
  participants 2, 4, 14, 18 held out
  (\texttt{configs/heldout\_subjects\_xu2025alljoined.json}),
  scored on their test-split windows
  (200 images × \textasciitilde80
  repetitions, 64,895 windows,
  \textasciitilde800 queries, SE ≈ 1.6
  points); training on the other 16
  participants (1,068,339 windows).
\item
  \textbf{First results (round 4):} A1
  (R6 recipe) \textbf{0.2775} (a
  one-epoch peak; neighbours ≈ 0.265),
  A2 (+ \texttt{window\_norm}) 0.2612
  (stopped early at 14). Inconclusive at
  this noise level; re-test with
  smoothed early stopping (§10.3).
\item
  \textbf{All 32 Alljoined channels
  exist in THINGS-EEG2\textquotesingle s
  63}, so THINGS-EEG2 can be reduced to
  the same montage:
  \texttt{-\/-extra-cache\ cache/windows/things\_eeg2}
  trains on Alljoined\textquotesingle s
  16 participants +
  THINGS-EEG2\textquotesingle s 10
  (train + val only), 1,736,739 windows
  (A4).
\end{itemize}

\subsubsection{10.3 New training options
(all off by default; defaults train
bit-identically)}\label{new-training-options-all-off-by-default-defaults-train-bit-identically}

{\def\LTcaptype{none} % do not increment counter
\begin{longtable}[]{@{}
  >{\raggedright\arraybackslash}p{(\linewidth - 4\tabcolsep) * \real{0.1379}}
  >{\raggedright\arraybackslash}p{(\linewidth - 4\tabcolsep) * \real{0.4310}}
  >{\raggedright\arraybackslash}p{(\linewidth - 4\tabcolsep) * \real{0.4310}}@{}}
\toprule\noalign{}
\begin{minipage}[b]{\linewidth}\raggedright
Option
\end{minipage} &
\begin{minipage}[b]{\linewidth}\raggedright
What it does
\end{minipage} &
\begin{minipage}[b]{\linewidth}\raggedright
Source
\end{minipage} \\
\midrule\noalign{}
\endhead
\bottomrule\noalign{}
\endlastfoot
\texttt{window\_norm} & per-window
channel z-scoring in \texttt{forward}
(also at inference; a marker buffer in
the weights switches it back on) &
standard EEG normalisation;
REVE\textquotesingle s preprocessing \\
\texttt{adversarial\_weight} &
\textbf{domain confusion}: a
training-only classifier learns the
participant from detached embeddings,
the encoder is pushed to make its output
uniform. A gradient-reversal version
(DANN) first made the classifier
confidently \emph{wrong} (accuracy below
chance), keeping the information;
neither removed it for a fresh probe on
a toy problem, and R19 shows it hurts &
Ganin \& Lempitsky 2015 (DANN); Tzeng et
al. 2015 (domain confusion); Elazar \&
Goldberg 2018 (adversarial removal often
fails) \textbf{(unverified)} \\
\texttt{smooth=k} & early stopping and
best-epoch selection on the mean of the
last k held-out scores (only full
windows); next rounds use
\texttt{smooth=3\ patience=10} &
motivated by R24/A1/A2 stopping on noise
spikes \\
\texttt{-\/-train-subjects\ N} & train
on the first N remaining participants
(participant curve) & \\
\texttt{-\/-extra-cache\ DIR} & joint
training with another
dataset\textquotesingle s window cache,
channels mapped by name, subject ids
offset & \\
\end{longtable}
}

\subsubsection{10.4 REVE transfer, stage
1 (round
5)}\label{reve-transfer-stage-1-round-5}

REVE (El Ouahidi et al., 2025;
\texttt{brain-bzh/reve-base}, 69M
parameters, pretrained on EEG from many
people) makes one 512-d token per
channel for a 1 s window, placing
channels by name (any montage). Frozen
and \textbf{averaged over channels} (the
official probe), it reached only
\textasciitilde0.13.

\begin{itemize}
\tightlist
\item
  \textbf{Stage 1} keeps every
  channel\textquotesingle s token and
  trains an \textbf{attention-pooling
  head} (\texttt{AttentionPoolHead}: 8
  learned queries attend over the
  channel tokens, then an MLP to the
  1536-d embedding):
  \texttt{scripts/cache\_reve\_tokens.py}
  runs frozen REVE once (official
  preprocessing: 200 Hz, per-window
  z-score, clamp ±15; float16 tokens),
  \texttt{scripts/train\_on\_tokens.py}
  trains the head with
  \texttt{Solver.fit} unchanged,
  \texttt{submission/reve\_submission.py}
  is the Codabench solver (the worker
  image installs braindecode and stages
  \texttt{brain-bzh/reve-base}; only the
  head is shipped). The exported
  submission on raw windows matches the
  head on cached tokens (cosine
  1.000000).
\item
  \textbf{Stage 2} (only if stage 1 is
  competitive): fine-tune
  REVE\textquotesingle s top layers on
  cached intermediate tokens. Full
  fine-tuning costs \textasciitilde2 h
  per Alljoined epoch on the Mac: too
  slow.
\end{itemize}

\subsubsection{10.5 Where this leaves
the final
recipe}\label{where-this-leaves-the-final-recipe}

In: per-window normalisation, wider
convolutions (if the slowdown is
affordable), more participants (joint
THINGS-EEG2 + Alljoined), smoothed early
stopping. Out: participant-adversarial
training, augmentation. Open: spatial
attention (only needed if the
final\textquotesingle s montage differs
from Alljoined\textquotesingle s), REVE
(round 5), repetition averaging (re-test
with smoothed early stopping).

\end{document}
